\documentclass[UTF8,a4paper,zihao=-4,fontset=none]{ctexart}
\usepackage[margin=25mm,headheight=16pt]{geometry}
\usepackage{fontspec}
\IfFontExistsTF{SimSun}{
  \setCJKmainfont{SimSun}[BoldFont=SimHei,ItalicFont=KaiTi]
  \setCJKsansfont{SimHei}\setCJKmonofont{FangSong}
}{
  \setCJKmainfont{FandolSong-Regular}[BoldFont=FandolSong-Bold,ItalicFont=FandolKai-Regular]
  \setCJKsansfont{FandolHei-Regular}\setCJKmonofont{FandolFang-Regular}
}
\IfFontExistsTF{Times New Roman}{\setmainfont{Times New Roman}}{\setmainfont{TeX Gyre Termes}}
\IfFontExistsTF{Arial}{\setsansfont{Arial}}{\setsansfont{TeX Gyre Heros}}
\IfFontExistsTF{Consolas}{\setmonofont[Scale=0.85]{Consolas}}{\setmonofont[Scale=0.85]{Latin Modern Mono}}
\usepackage{amsmath,amssymb,amsthm,bm,mathtools}
\usepackage{booktabs,longtable,tabularx,array,multirow}
\usepackage{graphicx,xcolor,float,caption,subcaption}
\usepackage[section]{placeins}
\usepackage{tikz,pgfplots}
\usetikzlibrary{arrows.meta,positioning,fit,calc}
\usepgfplotslibrary{groupplots}
\pgfplotsset{compat=1.18}
\usepackage[ruled,vlined,linesnumbered]{algorithm2e}
\usepackage{enumitem,fancyhdr,lastpage,listings}
\usepackage[numbers,sort&compress]{natbib}
\usepackage[colorlinks=true,linkcolor=navy,citecolor=teal,urlcolor=teal,breaklinks=true]{hyperref}
\usepackage{bookmark}
\definecolor{navy}{HTML}{17324D}
\definecolor{teal}{HTML}{087F8C}
\definecolor{muted}{HTML}{586579}
\definecolor{paleblue}{HTML}{EDF4F7}
\definecolor{orange}{HTML}{C47722}
\definecolor{brick}{HTML}{B24B4B}
\hypersetup{pdftitle={CrowdLocateLab：面向异步众包定位的科研智能体},pdfauthor={课程研究项目},pdfsubject={实际大模型智能体运行、稳健估计、几何可辨识性与分割共形校准}}
\ctexset{section={format=\Large\bfseries\color{navy}},subsection={format=\large\bfseries\color{navy}},subsubsection={format=\normalsize\bfseries}}
\newtheorem{proposition}{命题}[section]
\newtheorem{theorem}{定理}[section]
\newtheorem{lemma}{引理}[section]
\theoremstyle{remark}
\newtheorem{remark}{说明}[section]
\newtheorem{definition}{定义}[section]
\renewcommand{\proofname}{证明}
\renewcommand{\tablename}{表}
\renewcommand{\figurename}{图}
\SetAlgorithmName{算法}{算法}{算法目录}
\captionsetup{font=small,labelfont=bf,labelsep=quad}
\setlist{nosep,leftmargin=2em,topsep=4pt}
\setlength{\parskip}{3pt}
\linespread{1.16}
\setlength{\emergencystretch}{2em}
\pagestyle{fancy}
\fancyhf{}
\fancyhead[L]{\small\sffamily\color{muted}CrowdLocateLab}
\fancyhead[R]{\small\color{muted}科研智能体开发与应用}
\fancyfoot[C]{\small\thepage}
\renewcommand{\headrulewidth}{0.3pt}
\newcolumntype{Y}{>{\raggedright\arraybackslash}X}
\lstset{basicstyle=\ttfamily\footnotesize,breaklines=true,columns=fullflexible,keepspaces=true,frame=single,rulecolor=\color{black!15},backgroundcolor=\color{black!2},showstringspaces=false}
\pgfplotsset{every axis/.append style={font=\small,tick label style={font=\footnotesize},label style={font=\small},legend style={font=\footnotesize,draw=none},axis line style={black!35},grid style={black!10},major grid style={black!10}}}

\begin{document}
\begin{titlepage}
\thispagestyle{empty}
\setlength{\parindent}{0pt}
\vspace*{14mm}
{\sffamily\large\color{teal}AI FOR MATHEMATICS\hfill 课程研究项目\par}
\vspace{18mm}
{\Huge\bfseries\color{navy}CrowdLocateLab\par}
\vspace{6mm}
{\LARGE\bfseries 面向异步众包定位的\par 科研智能体开发与实证研究\par}
\vspace{8mm}
{\large\color{muted}从问题建模、假设修订到独立测试与误差校准\par}
\vspace{12mm}
\noindent\color{teal}\rule{\textwidth}{1.1pt}\color{black}
\vspace{6mm}
\begin{tabularx}{\textwidth}{@{}p{27mm}Y@{}}
研究来源 & 《HarmonyOS难题发布》第七期：基于非线性优化的众包离线高精度位置聚合难题\\[4mm]
研究对象 & 污染距离观测与运动模型失配下的当前位置估计\\[4mm]
系统形式 & 实际大模型规划器、研究工具、独立审阅调用、可恢复记忆与可核查轨迹\\[4mm]
交付内容 & Python源代码、真实运行记录、合成数据、独立评测、LaTeX源码与本文PDF
\end{tabularx}
\vfill
{\color{muted}研究版本 3.0\quad /\quad 2026年9月25日\par}
\vspace{3mm}
{\small 姓名、学号及课程信息请在提交前补充。\par}
\end{titlepage}
\pagenumbering{Roman}
\section*{摘要}
\addcontentsline{toc}{section}{摘要}
本文面向课程给出的众包离线高精度位置聚合开放问题，开发了使用实际大语言模型的科研智能体CrowdLocateLab。系统将问题阅读、可检验假设、组合方法注册、数值实验、失败诊断、独立审阅调用、数学检查和方法冻结串联起来，保留原始模型输出、工具输入输出及可恢复状态。两轮研究合计完成25个研究动作、27次实际模型调用和12次实验或消融调用；第二轮以第一轮失败为已知先验，再用新数据评价修订。

研究关注污染测距与历史观测下的当前位置估计，实现静态与恒速模型、距离与对数残差、稳健损失、截尾、子集初始化及观测交叉验证等可组合组件。理论部分给出自加权标量偏差、镜像不可辨识性、局部误差界，并推导镜像世界下接受率与错误率的限制。校准部分对比全局与可观测选模分组的分割共形区域，明确有限样本秩和保证条件。各数学结论的适用范围与实证结果分开陈述。

第一轮主测试显示总体改善，但静态退化、稀疏模型排除和子群欠覆盖促成第二轮修订。第二轮在720例新主测试上的平均误差为3.920米，相比WLS降低40.99\%；相对第一轮冻结方法的改善为0.145米，配对95\%区间$[-0.043,0.334]$跨零。8份观测的运动边界中，误差由10.788米降至6.217米，但仍高于恒速稳健基线的5.107米。静态正常误差相对WLS退化29.86\%，预设10\%容忍目标未达到。分组区域在1200例独立覆盖检查中，90\%与95\%名义水平的实际覆盖分别为88.17\%与95.00\%；分组没有消除真实运动异常场景的欠覆盖。

两轮冻结后评测及敏感性分析合计完成22052次定位计算，129项软件测试通过。交付包括完整源码、真实运行证据、逐例数据、只读研究浏览器、新观测定位入口及LaTeX源码。合成研究只回答原工业问题的一个可计算子问题，未完成完整工业题验收，也不将已有统计方法的组合称为原创通用定理。

\vspace{4mm}
\noindent\textbf{关键词：}科研智能体；异步定位；稳健估计；模型选择；可辨识性；分组共形预测

\clearpage
\tableofcontents
\clearpage
\pagenumbering{arabic}
\section{研究背景与问题定义}
\label{sec:background}

\subsection{从工业开放问题到可计算研究对象}

众包离线定位利用附近设备接收到的无线广播，为暂时无法联网的目标提供位置估计。课程参考材料《HarmonyOS难题发布》第七期第1页将这一任务表述为“基于非线性优化的众包离线高精度位置聚合难题”\cite{huaweiHarmonyChallenge7}：目标持续发送蓝牙广播，协助设备上报自身定位与接收信号强度；系统需要从时间、定位精度和可信程度不同的报告中估计目标当前位置。困难并不只来自测距噪声。报告者坐标本身存在误差，历史报告可能对应不同的目标位置，少量严重错误上报还可能破坏聚类和非线性拟合。

原材料给出的静态模型为
\begin{equation}
  \widehat x=\operatorname*{arg\,min}_{x\in\mathbb R^2}
  \sum_{i=1}^{N}\omega_i\bigl(f(r_i)-\|x-l_i\|_2\bigr)^2,
  \label{eq:original-problem}
\end{equation}
其中 $l_i$ 为报告者位置，$r_i$ 为接收信号强度，$f(r_i)$ 为反演距离，$\omega_i$ 由定位时间、方式、精度和置信度等信息决定。材料提出，超过10个协助设备时CEP90应在30米以内，并要求在标准场景不劣化的条件下将移动目标检测精度提升20\%。后一指标没有给出具体统计定义；材料也未提供原系统实现和真实验收数据。因此，不能用某个合成基准上的平均误差下降替代这项工业要求。

本文将研究对象限定为\emph{受污染、异步、二维距离观测下的当前位置估计及其可信度评估}。报告者可移动，但每条上报包含该次观测时的位置；目标在一个短时间窗口内以静态或恒速模型近似。输入是上报坐标、时间、距离及标定尺度，输出是当前位置、求解诊断和可选的预测区域。无线协议、实际硬件标定、楼层判断、地图约束和大规模产品部署不在本次实现范围内。这种限定保留了原题的非线性估计、运动失配、污染和几何退化等核心数学困难，同时允许在本地CPU上重复实验。

\subsection{相关研究及本文的技术位置}

接收信号强度与距离的关系通常在对数域建模。Patwari等讨论了RSS与到达时间观测下的位置估计、统计模型和误差界\cite{patwari2003relative}。当分贝域噪声近似高斯时，反演距离具有乘性误差，因此距离域的平方残差与对数残差对应不同的统计假设。报告者坐标误差、非视距传播和设备标定差异还会进一步改变有效噪声。本项目借助这一背景构建可控模拟，并不将其视为对真实BLE环境的充分刻画。

稳健估计为少量错误观测提供了已有工具。Soares与Gomes将Huber损失用于受污染的测距定位，并研究相应非凸问题的凸下界和分布式算法\cite{soares2016badapples}。Fischler与Bolles提出的RANSAC则通过子集拟合与全体一致性评价寻找不受大量粗差支配的模型\cite{fischler1981random}。本文采用的非线性稳健损失、子集初始化和截尾重拟合受这些思想启发，但数值实现没有继承相关文献中另一种算法的全局最优或收敛保证。

选择静态还是运动模型是一项预测问题。交叉验证提供了区分训练拟合和观测外预测的基本方法\cite{stone1974crossvalidatory}；它本身不会消除小样本、相依观测和模型失配。定位唯一性还受到几何信息限制，局部优化是否收敛与观测是否能区分两个不同位置是两回事\cite{eren2004rigidity}。对于不确定性，split conformal能够在交换性和方法冻结等条件下提供边际预测覆盖\cite{lei2018distributionfree}，但不能直接保证任一具体困难样本或任意拒绝后子集的条件覆盖。

由此，本文并不把某个现成损失函数称为新的数学发现。可检验的研究内容是：这些组件在同一个异步定位问题中如何相互作用，哪些失败可由精确反例解释，以及工具使用型科研智能体能否根据真实证据作出有依据的修订。

\subsection{先导实验暴露的失败}

第一版原型采用固定候选库中的经验规划，尚未调用大语言模型。其独立合成测试结果见表\ref{tab:pilot-failure}。这些数字属于历史先导实验，不属于本文后续新版系统的最终测试。

\begin{table}[tbp]
  \centering
  \caption{第一版先导实验的历史结果：160个合成测试案例。误差越低越好。}
  \label{tab:pilot-failure}
  \begin{tabular}{lrr}
    \toprule
    方法 & 平均误差（m） & CEP90（m）\\
    \midrule
    静态WLS & 6.4840 & 14.3271\\
    经验规划所选配置 & 4.3632 & 12.4707\\
    同预算随机搜索所选配置 & 3.6839 & 10.2265\\
    固定恒速稳健模型 & 4.1668 & 9.5568\\
    \bottomrule
  \end{tabular}
\end{table}

经验规划虽然相对静态WLS改善总体均值，却没有优于随机搜索或固定恒速稳健模型。在移动异常场景中，启发式模型评分将40例中的37例选择为静态，该场景均值为11.2833米，略高于WLS的11.2000米；静态正常场景也从0.8743米退化到0.9113米。这表明，总体均值下降既不能证明搜索策略更有效，也不意味着实现了原题的标准场景非劣要求。

这些观察引出三个待检验解释：由观测距离自身构造尺度可能改变误差权重；污染下的训练拟合评分可能误判运动；有限开发案例上的选择可能发生泛化失败。它们不是已经确定的因果结论。只有保持其余组件不变的消融，以及独立案例上的验证，才能区分各项解释。因为这些历史测试结果已被研究者看到，它们可以作为诊断材料，不能再次被称为新版方案完全未见的最终测试。

\subsection{研究问题与贡献边界}

本文围绕以下相互关联的问题组织研究。
\begin{enumerate}
  \item 对数残差及其正确中心化，能否改善乘性误差下的拟合，并减少观测依赖权重导致的系统偏差？
  \item 在相同求解组件下，观测交叉验证能否比训练残差驱动的门控更好地权衡静态稳定性与运动解释能力？
  \item 局部几何诊断、近似预测尺度与独立校准，能够支持什么范围的可信度结论；哪些歧义无法仅靠数值优化消除？
  \item 一个实际调用语言模型、实验与核验工具的科研控制器，是否形成可复盘的假设、行动、观察与修订链条？
\end{enumerate}

数学贡献采用明确假设下的项目内推导：标量相对平方风险的偏差反例、恒速镜像歧义及不可区分性下界、包含曲率和邻域条件的局部误差界。它们用于解释与约束实验，不作理论首创声明。系统贡献是独立实现的工具调用运行时、可组合定位算子、证据与状态记录，以及真实推理和日志回放的区分。

架构参考来自Creative-Intelligence所强调的研究、验证和持久状态分离思想\footnote{参考项目：\texttt{github.com/DeepMathLLM/Creative-Intelligence}，查阅日期为2026年9月25日。}。本项目不是该仓库的fork，没有复制其源码，也不以运行Moonshine作为前提。本文对智能体实际行为的描述以保存的模型请求、响应和工具证据为依据；模型给出的评论既不是独立实验复现，也不是数学证明。

\section{数学建模、方法组合与可验证结论}
\label{sec:theory}

\subsection{异步观测、对数残差与解析导数}

以下以 $y_i>0$ 表示反演距离，以 $a_i\in\mathbb R^2$ 表示上报坐标，当前时刻为0，历史时刻满足 $t_i\leq0$。静态模型的参数是 $x_0$；恒速模型的参数是 $\theta=(x_0,v)$，对应距离
\begin{equation}
 d_i(\theta)=\|x_0+v t_i-a_i\|_2.
 \label{eq:trajectory-range}
\end{equation}
真实位置与污染标签不进入定位器。上报坐标可带误差，因此式\eqref{eq:trajectory-range}中的 $a_i$ 不是默认无误差的物理基站；相关误差进入观测模型失配。

对数距离传播模型可写为
\begin{equation}
 p_i=P_0-10n_p\log_{10}(d_i/d_0)+\varepsilon_i,
 \qquad \varepsilon_i\sim N(0,\sigma_p^2),
 \label{eq:rssi-model}
\end{equation}
其中 $p_i$ 是RSSI，$n_p$ 为路径损耗系数。已知参数的直接反演导致对数距离上的高斯噪声\cite{patwari2003relative}。一般记
\begin{equation}
 \log y_i=\log d_i+\mu_i+e_i,
 \qquad e_i\sim N(0,s_i^2).
 \label{eq:log-observation}
\end{equation}
直接反演的典型约定是 $\mu_i=0$；若模拟器采用均值为1的距离乘子，则 $\mu_i=-s_i^2/2$。后者需要在残差中作负号修正：$\log d_i-\log y_i-s_i^2/2$。距离无偏与对数中心为零不是同一条件。

实现允许距离残差与对数残差，分别为
\begin{align}
 r_i^{\mathrm{range}}(\theta)&=(d_i(\theta)-y_i)/q_i,\\
 r_i^{\log}(\theta)&=
 \frac{\log\max(d_i(\theta),d_{\min})-
       \log\max(y_i,d_{\min})-b}{q_i},
 \label{eq:implemented-log-residual}
\end{align}
其中 $d_{\min}=0.1$ 米，修正开启时 $b=s_0^2/2$，关闭时 $b=0$；$s_0$ 对应配置中的标定对数噪声。对于距离残差，报告尺度为 $q_i=\sigma_i$；对于对数残差，实现使用
\begin{equation}
 q_i=\max\left\{10^{-5},\sqrt{\log\left(1+
       \left(\frac{\sigma_i}{\max(y_i,0.1)}\right)^2\right)}\right\}.
 \label{eq:implemented-log-scale}
\end{equation}
这是由相对误差构造的观测尺度近似，不是已知的真实对数方差。另一选项将 $q_i$ 固定为该拟合数据内尺度的中位数。偏差修正参数 $s_0$ 与式\eqref{eq:implemented-log-scale}中的有效尺度不同；当存在坐标误差、加性误差或标定失配时，不能把该实现称为精确最大似然。

令 $u_i=(x_0+vt_i-a_i)/d_i$。在 $d_i>0$ 处，距离的参数导数为
\begin{equation}
 D_\theta d_i=[u_i^{\mathsf T},\ t_i u_i^{\mathsf T}],
 \qquad
 D_\theta\log d_i=d_i^{-1}[u_i^{\mathsf T},\ t_i u_i^{\mathsf T}].
 \label{eq:range-jacobian}
\end{equation}
静态模型只保留位置列，标准化后再除以 $q_i$。实现对硬截断区 $d_i\leq d_{\min}$ 使用零导数；截断边界不可微，因而随后理论命题的光滑邻域必须远离该边界。解析导数用于数值求解，数值差分可以检验其局部实现，却不能证明优化器找到全局最小值。

\subsection{稳健拟合、截尾与子集初始化}

使用SciPy的平方残差接口记损失为 $\rho(z)$，$z\geq0$。当前实现的线性、soft-L1、Huber和Cauchy损失依次为
\begin{equation}
 z,\qquad 2(\sqrt{1+z}-1),\qquad
 \begin{cases}z,&z\leq1,\\2\sqrt z-1,&z>1,\end{cases}
 \qquad \log(1+z).
 \label{eq:loss-family}
\end{equation}
Huber测距定位已有相关研究\cite{soares2016badapples}。所有这些几何目标仍可能非凸；更强的降权还可能牺牲有效运动信息。

给定截尾比例 $\gamma$ 和参数维数 $p$，保留数为
\begin{equation}
 h=\min\{n,\max\{p+1,\lceil(1-\gamma)n\rceil\}\}.
\end{equation}
令 $I_h(\theta)$ 是绝对标准化残差最小的 $h$ 个观测，候选评分为
\begin{equation}
 L(\theta)=\sum_{i\in I_h(\theta)}
 C^2\rho\left(\frac{r_i(\theta)^2}{C^2}\right)
 +\lambda\|v\|_2^2,
 \label{eq:trimmed-robust-objective}
\end{equation}
静态模型没有速度惩罚。正则项始终保持二次，不随观测稳健损失一起降权。求解过程从多个观测导出的起点开始，先拟合再更新保留集合；启用截尾时最多执行有限次重拟合，最后按式\eqref{eq:trimmed-robust-objective}选择候选。有限轮交替不等于截尾目标的全局求解，也不保证返回值已经达到保留集合不变的固定点。

可选的子集初始化受抽样一致思想启发\cite{fischler1981random}，但当前实现不是经典硬阈值RANSAC。每次均匀抽取
\begin{equation}
 m=\min\{n,\max\{p+2,\lceil n/2\rceil\}\}
\end{equation}
个观测，在子集上拟合；各结果以全体观测上的截尾稳健目标评分，最佳结果作为一个额外起点，再进行正常重拟合。随机性由观测字节内容及声明种子共同确定，不读取真实场景标签。它没有按时间或几何预筛选子集，因此一次抽到内点不等于抽到可辨识且条件良好的子问题。

若恰有 $h_0$ 个内点、每次独立均匀抽取大小为 $m$ 的子集，全内点概率为 $\binom{h_0}{m}/\binom{n}{m}$。多次抽样的成功事件只说明出现过全内点子集，不证明该子集拟合成功或后续评分能识别正确轨迹。截尾数、子集数、起点数和迭代上限均影响计算量；公平比较不能只数最终登记了多少个方法。

\subsection{以留出观测比较静态和恒速模型}

模型选择采用案例内的观测交叉验证，而不是直接比较两个模型的训练拟合损失。交叉验证是预测评估方法\cite{stone1974crossvalidatory}；本实现首先按时间排序，再交错分成 $K$ 折，使每个观测恰好作为验证点一次。每折在剩余观测上分别拟合静态模型 $M_0$ 和恒速模型 $M_1$，折内使用一个起点，并限制子集初始化次数；最终选定模型再以完整声明预算在全部观测上重拟合。因此，这里验证的是计算预算缩减的代理拟合过程，不是对完整最终求解过程逐字等同的交叉验证。

验证评分不使用拟合时的截尾，也不跟随候选优化损失改变。它始终对全部留出标准化残差使用尺度为1的soft-L1：
\begin{equation}
 \widehat R_M=\frac1n\sum_{k=1}^{K}\sum_{i\in V_k}
 2\left(\sqrt{1+r_i(\widehat\theta_{M,-k})^2}-1\right).
 \label{eq:cv-score}
\end{equation}
若采用常数尺度，验证折沿用训练折的常数；报告尺度模式则从该条验证观测的公开距离与尺度计算。满足训练样本数与时间跨度要求、各折恒速拟合未回退、数据Jacobian满列秩且没有活跃数值边界时，才允许比较
\begin{equation}
 \widehat R_{M_1}<(1-\delta)\widehat R_{M_0},
 \qquad 0\leq\delta<1.
 \label{eq:cv-gate}
\end{equation}
否则选静态模型。这里 $\delta$ 是相对收益门槛，不是显著性水平；实现没有在这一门控中额外施加条件数阈值。

训练和验证分离降低了直接奖励模型自由度的倾向，但式\eqref{eq:cv-gate}不是运动假设检验。相依上报、污染验证点、标定误差和加速轨迹仍会使其失效。交错折主要检验窗口内部的预测，不等价于从早期观测外推到当前时刻；折间训练集重叠，也不能把各折当成独立样本计算常规显著性。最终位置误差的改进需要独立案例验证。

\subsection{三个精确结论：偏差、歧义与局部稳定性}

\begin{proposition}[由观测自身决定尺度的标量偏差]
\label{prop:self-weight-bias}
设 $d>0$，$Y=d\exp(\varepsilon+\mu)$，$\varepsilon\sim N(0,s^2)$ 且 $s>0$。在 $q>0$ 上，相对平方风险
\begin{equation}
 R(q)=\mathbb E\left[\frac{(q-Y)^2}{Y^2}\right]
\end{equation}
的唯一最小点为 $q_*=d\exp(\mu-3s^2/2)$。正确中心化的对数平方风险 $\mathbb E[(\log q-\log Y+\mu)^2]$ 则在 $q=d$ 处唯一最小。特别地，均值为1的乘子对应 $\mu=-s^2/2$，此时 $q_*=d e^{-2s^2}$。
\end{proposition}

\begin{proof}
展开风险得 $R(q)=q^2\mathbb E[Y^{-2}]-2q\mathbb E[Y^{-1}]+1$。其二阶导数严格为正，唯一驻点为 $\mathbb E[Y^{-1}]/\mathbb E[Y^{-2}]$。由高斯指数矩，
\[
 \mathbb E[Y^{-1}]=d^{-1}e^{-\mu+s^2/2},\qquad
 \mathbb E[Y^{-2}]=d^{-2}e^{-2\mu+2s^2}.
\]
相除即得第一式。另一方面，$\log Y=\log d+\mu+\varepsilon$，故中心化的对数风险等于 $(\log q-\log d)^2+s^2$，其唯一最小点为 $d$。
\end{proof}

\begin{remark}
该结论给出观测依赖权重的一个精确简化反例，并非完整定位器的偏差公式。式\eqref{eq:implemented-log-scale}还含加性尺度近似；实际多维目标具有坐标误差、非线性几何、正则项和局部极小值。不能由命题\ref{prop:self-weight-bias}推断对数方法在所有场景均优，也不能把它当成第一版全部失败的已证实原因。
\end{remark}

\begin{proposition}[动态镜像歧义与不可区分性下界]
\label{prop:mirror}
设全部已知报告者坐标为 $a_i=(b_i,0)$，$R=\operatorname{diag}(1,-1)$。两条恒速轨迹 $(x_0,v)$ 与 $(Rx_0,Rv)$ 在任意观测时刻产生完全相同的精确距离。若其当前位置间距为 $D>0$，任意确定点估计至少对其中一个真值产生 $D/2$ 的欧氏误差。若加入同一仅依赖距离的噪声机制，则任意随机估计器至少在其中一个世界具有不小于 $D/2$ 的期望误差。
\end{proposition}

\begin{proof}
由 $Ra_i=a_i$ 及反射保持范数，
\[
 \|Rx_0+Rv t_i-a_i\|
 =\|R(x_0+v t_i-a_i)\|=\|x_0+v t_i-a_i\|.
\]
对于同一观测得到的点估计 $\widehat x$，三角不等式给出 $D\leq\|x_0-\widehat x\|+\|Rx_0-\widehat x\|$，所以至少一个误差不小于 $D/2$。两个世界加入题设噪声后具有相同观测分布；再对同一观测与估计器随机性积分，两个期望误差之和仍至少为 $D$，从而得到最后一个结论。
\end{proof}

局部满列秩不能排除这种全局歧义。例如，取 $x_0=(1.5,3)$、$v=(0.2,0.1)$、$b=(-7,-2,5,-5,7,1)$ 和 $t=(-8,-6,-4,-2,-1,0)$，距离Jacobian可以为四列满秩，而镜像轨迹仍给出相同距离。正的逐行标准化及对数变换不改变该秩。这里的反例由等距关系直接证明；网络刚性理论提供了更一般的唯一定位背景\cite{eren2004rigidity}。增加同轴观测或更换稳健损失不会产生原本缺失的两侧信息。

\begin{figure}[htbp]
\centering
\begin{tikzpicture}[x=0.75cm,y=0.70cm,font=\small]
\draw[->,black!35] (-8,0)--(8.3,0) node[right] {$x$};
\draw[->,black!35] (0,-4)--(0,4) node[above] {$y$};
\foreach \x in {-7,-2,5,-5,7,1} {\fill[navy] (\x,0) circle (2.3pt);}
\draw[-{Stealth[length=2.5mm]},thick,teal] (-.1,2.2)--(1.5,3);
\draw[-{Stealth[length=2.5mm]},thick,orange] (-.1,-2.2)--(1.5,-3);
\fill[teal] (1.5,3) circle (2.5pt) node[above right] {$x_0=(1.5,3)$};
\fill[orange] (1.5,-3) circle (2.5pt) node[below right] {$Rx_0=(1.5,-3)$};
\draw[dashed,navy!50] (-2,0)--(1.5,3);
\draw[dashed,navy!50] (-2,0)--(1.5,-3);
\node[align=left,text=muted] at (-5,2.7) {同轴报告者\\观测距离无法区分两侧};
\draw[<->,brick] (6,3)--(6,-3) node[midway,right,fill=white] {$D=6\,\mathrm m$};
\end{tikzpicture}
\caption{动态镜像反例。报告者横坐标与观测时刻取命题中的六组数值；上下两条轨迹对应各时刻的测距完全相同。两条轨迹的当前位置相隔6米，任何只看这些观测的点估计至少在一个世界中误差不小于3米。}
\label{fig:mirror}
\end{figure}


\begin{proposition}[含曲率条件的局部非线性误差界]
\label{prop:local-bound}
取固定时间尺度 $T>0$，令 $\bar\theta=(x_0,u)$、$u=Tv$、$\tau_i=t_i/T$，并定义
\[
 A_i=[I_2\ \tau_i I_2],\qquad
 F_i(\bar\theta)=s_i^{-1}\log\|A_i\bar\theta-a_i\|,
 \quad s_i>0.
\]
设在闭球 $B(\bar\theta_*,R_0)$ 内所有距离至少为 $r_{\min}>0$，$J_*=DF(\bar\theta_*)$ 的最小奇异值为 $\sigma_*>0$。令
\begin{equation}
 L=\left[\sum_i\left(\frac{1+\tau_i^2}{s_i r_{\min}^2}\right)^2\right]^{1/2}.
\end{equation}
对球内的 $\bar\theta_*+h$，有
\begin{equation}
 \|F(\bar\theta_*+h)-F(\bar\theta_*)\|
 \geq\sigma_*\|h\|-\tfrac12L\|h\|^2.
 \label{eq:nonlinear-local-lower}
\end{equation}
进一步，若观测向量 $b=F(\bar\theta_*)+e$ 满足 $\|e\|\leq\eta$，估计值位于上述球内、$\|\widehat{\bar\theta}-\bar\theta_*\|\leq\sigma_*/L$，且 $\|F(\widehat{\bar\theta})-b\|\leq\beta$，则
\begin{equation}
 \|\widehat x_0-x_{0,*}\|
 \leq\|\widehat{\bar\theta}-\bar\theta_*\|
 \leq\frac{2(\eta+\beta)}{\sigma_*}.
 \label{eq:conditional-position-bound}
\end{equation}
\end{proposition}

\begin{proof}
对于 $z\ne0$，令 $q=z/\|z\|$，则
\[
 \nabla^2\log\|z\|=(I-2qq^{\mathsf T})/\|z\|^2,
 \qquad \|\nabla^2\log\|z\|\|_2=\|z\|^{-2}.
\]
又 $\|A_i\|_2^2=1+\tau_i^2$，每个 $F_i$ 的Hessian在该球内至多为 $(1+\tau_i^2)/(s_i r_{\min}^2)$。Taylor余项组成的向量范数不超过 $L\|h\|^2/2$。结合 $\|J_*h\|\geq\sigma_*\|h\|$，得到式\eqref{eq:nonlinear-local-lower}。在 $\|h\|\leq\sigma_*/L$ 时，下界至少为 $\sigma_*\|h\|/2$；另一方面，三角不等式将映射差上界控制为 $\eta+\beta$，从而得到式\eqref{eq:conditional-position-bound}。
\end{proof}

若另有 $\beta\leq\eta$，上界可化为 $4\eta/\sigma_*$。但鲁棒目标的小值不自动意味着未截尾二范数残差满足这一条件。真值邻域、噪声界与真值处Jacobian通常也未知，所以该命题是有条件的局部结论，不是把拟合点奇异值代入即可得到的全局证书。它不排除命题\ref{prop:mirror}中的远处镜像，也不适用于跨越硬截断边界的邻域。时间重参数化用于使条件数与状态范数具有明确尺度；实现直接优化 $v$，若换成 $u$，同一速度惩罚应相应写成 $\lambda\|u\|^2/T^2$。

\subsection{从几何尺度到校准预测区域}

定位模块输出一个\emph{未校准}的几何尺度。记保留观测的数据Jacobian为 $J_I$，不包含速度正则行，$r$ 则包含全部观测的标准化残差。满列秩时，实现使用
\begin{equation}
 g(Z)=\operatorname{clip}_{[0.2,100]}\left[
 \sqrt{\operatorname{tr}\left\{(J_I^{\mathsf T}J_I)^\dagger_{xx}\right\}}
 \max\left\{1,\frac{\operatorname{median}_i|r_i|}{0.6745}\right\}\right],
\end{equation}
其中下标 $xx$ 表示伪逆的位置块；秩亏时返回上限100米。这个量只是观测驱动的代理量，不是Cram\'{e}r--Rao下界或95\%置信半径；固定模型下的局部线性近似也没有处理门控、不正确内点集合和全局多解。

一个独立校准层可以将其转化为案例级预测圆盘。方法和 $g$ 冻结后，取 $m$ 个未参与选择的校准案例，定义
\begin{equation}
 S_j=\frac{\|\widehat x(Z_j)-x_j\|_2}{g(Z_j)},\quad
 k=\lceil(m+1)(1-\alpha)\rceil,\quad q=S_{(k)},
 \label{eq:conformal-score}
\end{equation}
若 $k>m$，规定 $q=+\infty$。圆盘为
\begin{equation}
 C(Z)=\{x:\|x-\widehat x(Z)\|_2\leq qg(Z)\}.
\end{equation}
在校准与新案例可交换，且估计器和尺度函数已独立冻结的条件下，split conformal的秩论证给出
\begin{equation}
 \Pr\{x_{\mathrm{new}}\in C(Z_{\mathrm{new}})\}\geq1-\alpha.
\end{equation}
原因是新分数在 $m+1$ 个可交换分数中的秩对称；上述取整提供保守覆盖，平局不会降低覆盖下界。这是已有预测覆盖思想在二维欧氏非符合度上的直接构造\cite{lei2018distributionfree}，并非本项目新定理，也不是单个条件样本的覆盖保证。

校准单位应是完整定位案例，而非同一轨迹中的各条上报。按真实场景固定配额构造的分层样本，不自动满足合并分数的交换性；若部署分布或压力分布变化，亦不能无条件沿用上述保证。校准结果可以实证评价，但保证范围必须与实际抽样协议一致。

以大预测半径、数值边界或拟合分歧拒绝定位后，需要同时报告接受率、被接受案例误差和区域覆盖。总体的边际覆盖不自动成为拒绝后子集的条件覆盖。几何拒绝也可能遗漏局部满秩的镜像歧义。由此留下的研究问题是：能否在明确污染类、时间跨度、几何与计算预算条件下，兼顾静态非劣、运动收益和有意义的拒绝率？本文的命题给出了部分限制和局部工具，并未解决这一完整问题。

\section{工具使用型科研智能体的设计与审计}
\label{sec:agent-system}

\subsection{模型决策、数值行动与证据的分离}

系统将研究规划与定位求解分为两个不同层次。语言模型接收目标、方法登记表、近期实验观察、持久记忆和审阅意见，自行提出下一步研究行动；外层控制器解析并校验响应，调用登记的确定性工具，再把工具结果送回下一次模型请求。数值结果来自实际工具执行，不能由模型响应中的文字代替。算法\ref{alg:research-loop}给出这一闭环的结构。

用户提供的Creative-Intelligence项目\cite{creativeIntelligence}为研究流程与状态管理的分工提供了架构参考。本项目重新实现控制器、任务专用数学工具与评测协议，没有复制该仓库源码，也不把自身称为Moonshine系统的分支。

与固定候选列表的自动调参相比，本系统允许模型同时提出可证伪假设、解释预计发生的现象、登记一个已实现算子的组合，并根据失败结果修改方法。其自主性仍然有明确边界：模型不能在这个协议中任意生成并执行新的Python程序，也不能增添未经实现的求解器。所谓方法构造，是在静态、恒速与交叉验证门控、距离或对数残差、尺度形式、稳健损失、截尾和子集初始化等已实现机制之间组合。数学创新程度因而需要由具体研究成果判断，不能从“使用了语言模型”直接推出。

\begin{figure}[H]
\centering
\begin{tikzpicture}[box/.style={draw=navy!50,fill=paleblue,rounded corners=2pt,align=center,text width=37mm,minimum height=13mm,font=\small},arr/.style={-{Stealth[length=2mm]},draw=teal,thick},node distance=9mm and 12mm]
\node[box] (plan) {大模型规划器\\形成假设、选择动作};
\node[box,right=of plan] (control) {持久化控制器\\校验、预算、派发};
\node[box,right=of control] (tool) {研究工具与求解器\\实验、诊断、数学见证};
\node[box,below=of tool] (evidence) {真实实验结果\\逐例数据与诊断};
\node[box,below=of control] (memory) {记忆与审阅反馈\\保留失败，修订假设};
\node[box,below=of plan] (freeze) {方案冻结\\独立测试与校准};
\draw[arr] (plan)--(control);
\draw[arr] (control)--(tool);
\draw[arr] (tool)--(evidence);
\draw[arr] (evidence)--(memory);
\draw[arr] (memory.north)--++(0,4mm)-|([xshift=-4mm]plan.south);
\draw[arr] ([xshift=5mm]plan.south)--([xshift=5mm]freeze.north);
\end{tikzpicture}
\caption{智能体的研究闭环。动作来自实际模型响应，证据来自数值工具；冻结后才开放最终评测。}
\label{fig:architecture}
\end{figure}


\begin{table}[H]
  \centering
  \small
  \caption{运行时各组件的职责与证据边界。}
  \label{tab:agent-components}
  \begin{tabular}{p{0.18\linewidth}p{0.35\linewidth}p{0.36\linewidth}}
    \toprule
    组件 & 主要职责 & 输出和限制\\
    \midrule
    模型提供者 & 发出真实CLI请求，或显式回放既有响应 & 原请求、JSON响应、事件与用量；回放不算新推理\\
    研究控制器 & 校验动作、检查预算、派发工具、维护状态 & 事件链与状态快照；模型失败不切换成规则提案\\
    研究工具箱 & 登记方法、执行实验、诊断、消融、冻结 & 指标与逐例证据；开发接口不能选择测试或校准集\\
    定位模块 & 根据观测求解声明的方法组合 & 位置、诊断、工作量；不读取真实位置和污染标签\\
    审阅者调用 & 审查泄漏、对照、计算预算和结论范围 & 同模型的独立提示调用；不等于独立复现\\
    数学核验工具 & 执行指定反例及有限数值见证 & 数值支持或失败；不替代命题证明\\
    \bottomrule
  \end{tabular}
\end{table}

\begin{algorithm}[H]
  \small
  \DontPrintSemicolon
  \caption{带结构校验、证据记录和冻结边界的科研闭环}
  \label{alg:research-loop}
  \KwData{研究目标、登记工具、开发数据接口、计算预算与模型提供者}
  \KwResult{具有验证证据的冻结方法，或明确记录的停止状态}
  初始化或核验恢复状态、事件链和运行身份摘要\;
  \While{状态为运行中且尚未冻结}{
    检查人工停止标记及动作、实验、模型调用预算\;
    \If{达到预设的成功常规实验审阅间隔}{
      发起独立审阅者模型调用，保存其意见并加入记忆\;
    }
    将工具定义、登记方法、近期观察和持久记忆交给规划模型\;
    校验响应字段、动作枚举和参数JSON，保存原始模型响应\;
    \If{请求失败或结构非法且有限重试已耗尽}{
      记录提供者失败，不生成替代研究决策\;
      \Return{失败状态及已经保存的证据}\;
    }
    写入待执行工具标记和调用事件，再执行登记的研究工具\;
    保存工具结果或错误观察，更新证据索引与记忆\;
    \If{冻结动作通过工具箱的证据门槛}{
      保存方法及登记表摘要，关闭后续开发动作\;
    }
  }
  返回冻结状态或失败、预算耗尽、人工停止等实际状态\;
\end{algorithm}

\subsection{结构化决策与可组合工具}

规划器的响应包含公开决策摘要、研究假设、动作名称、参数JSON、预期观察和记忆备注。动作只能取自登记集合。控制器要求所有必填字段恰好存在、字段类型正确、动作合法且参数可解析为JSON对象，并拒绝非有限JSON数值。参数进入工具后，还由工具签名与定位配置验证再次约束。这里保存的是供研究者审阅的简洁决策理由，不是要求模型暴露内部思维链。

\texttt{read\_sources}返回研究问题与实现数学说明。来源文字属于证据材料，不拥有改变运行规则的权限。\texttt{register\_method}接收方法标识、假设、数学定义和声明式配置；同一运行中既有标识不能覆盖，因此修改后的方法具有新的版本标识。登记时用一个无关的小案例检查接口是否能运行，这只是接口探测，不计入研究表现。

\texttt{run\_experiment}在相同案例上评价指定方法，返回总体与分场景误差、尾部指标、工作量及证据路径。开发期只有发现集和验证集接口，验证调用还有独立次数限制。\texttt{diagnose}在已执行结果中读取失败案例、模型选择频率与求解诊断；其看到的真实误差来自离线评价端，而求解器只接触观测字段。这一区分允许智能体研究开发集上的失败，却不把未知案例真值作为部署输入。

\texttt{ablate}复制已登记配置并施加模型提出的组件修改，然后在共同的发现集上比较两者。该工具保存改动内容，但因果解释仍需检查是否只改变了待研究因素；若一次同时修改多个字段，就不能归因于单一机制。\texttt{verify\_math}执行预先实现的数值见证。它检验给定推导是否出现数值矛盾，不自动证明任意新命题。

\texttt{freeze}将选择的方法、科学结论、验证结果和登记表摘要写入冻结文件。工具箱检查候选是否具有验证证据，以及是否完成最低限度的实验、多方法、诊断和数学核验。通过后开发工具关闭，独立测试与校准位于研究控制器之外。冻结门槛保证存在可审查过程，不保证选中了最优方法；负面或无显著改善的结论同样可能是合法研究结果。

\subsection{审阅者、记忆与方法修订}

审阅者是与规划步骤分开的模型调用，使用专门提示审查数据泄漏、缺少对照、计算预算不公平、未经证明的创新声明和遗漏的反例。它使用相同的配置模型，而不是一个独立训练的第二模型；“独立”仅指调用上下文和职责分开。它可能与规划器共享盲点，其意见本身既不是证明，也不构成独立实验复现。

审阅意见与每步假设、预期观察、执行状态一起进入持久记忆；下一次规划上下文还包含方法登记表、最近的工具观察和剩余预算。模型可以据此提出新的组合、修改参数、安排消融或停止某条思路。历史结果不会被替换成一段只保留成功的最终叙述：完整调用和失败均留在事件记录中。当前控制器按成功的常规实验调用触发周期审阅；消融工具虽然消耗实验预算，但不能把它的存在误报为一次额外审阅者推理。

这种结构支持对研究过程的几个独立问题进行核查：提案是否先于对应结果；修订是否引用了当时已经存在的证据；某个失败是否导致后续行动改变；最终方法是否在测试访问之前冻结。它并不要求每轮都改善结果。若模型只是反复登记近似配置，日志同样会暴露这一点，不能因为系统具有“记忆”和“审阅者”组件就断言其研究有效。

\subsection{持久化、故障与审计强度}

真实请求保存提示文本、输出结构、响应JSON、CLI原始事件、标准错误和元数据。元数据记录提供者、请求摘要、耗时以及能够获得的用量。控制器另外维护追加式事件记录，每条事件包含序号、上一事件摘要和自身规范JSON摘要；状态快照采用临时文件写入后替换。此设计能够发现普通日志损坏、丢序或与保存摘要不一致的修改，但它不是数字签名或第三方时间戳，不能声称对能重写全部文件的操作者具有不可篡改保证。

恢复运行时，控制器核对研究目标、工具定义、配置、种子摘要、提供者和预算等运行身份信息，并验证事件链。工具执行之前先保存待执行标记；如果程序在工具完成与状态提交之间中断，重启会拒绝盲目重做这一副作用不确定的调用。该设计牺牲了一部分自动恢复便利，以避免重复登记、覆盖或把未完成结果当成已完成证据。

模型请求允许有限重试。超过限制、超时或提供者失败时，系统保存失败状态并停止相关流程，不悄悄退回UCB、随机提案、硬编码脚本或伪造模型响应。工具错误则成为下一轮可见的观察，由模型决定是否修正后再调用。总动作数、实验调用数、模型调用数和人工停止标记共同约束运行；预算耗尽本身不等同于完成研究，也不会授权尚未冻结的方案访问最终测试。

真实CLI调用请求只读沙箱，并在提示中禁止模型直接调用CLI内置的文件、命令或网络工具，由外层控制器执行研究工具。然而，提示不构成完全隔离的技术保证；提供者会保留原始CLI事件并标记意外工具行为。如果出现此类记录，就不能继续声称本次运行是完全无内置工具的纯模型接口。具体是否发生、模型版本和用量，应由本次运行证据说明。

\subsection{真实推理与显式回放}

回放模式顺序读取已经保存的模型响应，用于检查协议、工具调度和结果再计算。它在状态和请求审计中标记为回放，保留来源请求索引，并声明没有进行新的模型推理；既有响应耗尽即失败，不自动切换到在线调用或脚本答案。可选的严格模式要求提示摘要与原请求一致，非严格模式则记录不一致，不能把相近上下文下重用响应等同于模型在新证据下重新决策。

因此，回放成功只能支持“这组已保存决策能够被重新执行或检查”，不能单独证明语言模型今天仍会选择相同行动，也不能替代真实运行的模型事件。本文的实证部分将模型产生的新决策、确定性数值计算、数值反例验证与回放审计作为不同证据类别。这样的区分使研究深度能够由实际行为和结果判断，而不依赖未经验证的系统标签。

\section{第一轮独立评估协议与统计解释}
\label{sec:experimental-protocol}

\subsection{数据用途分离与冻结边界}

实验区分方法发现、开发验证、最终点估计比较以及预测区域校准。研究工具箱只提供发现集与验证集接口：前者用于提出和修订假设，后者允许有限次评价，因此也属于开发过程，不能称为最终无偏测试。方法冻结后，\texttt{benchmark\_v2\_1.py}读取冻结文件及登记表，写入协议、方法配置、源代码摘要和运行环境，然后生成独立评估案例。表\ref{tab:dataset-protocol}列出当前默认协议；运行时参数如有改变，应以该次保存的\texttt{protocol.json}为准。

评估器在研究尚未冻结、任何正式测试案例尚未生成时接受一次预先代码审查修订：当最终规格为旧基线简称时，不生成新组合方法的非法组件消融；正常组合方法分支完全不变。正式首次测试统一采用\texttt{benchmark\_v2\_1.py}，不比较两个版本后择优。原\texttt{benchmark.py}原字节保留为初始化版本与共享校验、统计、校准库；新入口直接复用这些函数。原快照、精确差异、修订时间和哈希保存于\texttt{research/}，新评估协议同时记录原研究manifest及新评估器版本。研究工具、数值求解器、模型轨迹、全部种子与指标均未修改。

\begin{table}[tbp]
  \centering
  \caption{默认数据协议。表中为设计参数，不是性能结果。}
  \label{tab:dataset-protocol}
  \begin{tabular}{p{0.16\linewidth}p{0.24\linewidth}p{0.29\linewidth}r}
    \toprule
    数据用途 & 场景与抽样方式 & 生成种子 & 案例数\\
    \midrule
    发现集 & 四个核心场景，各16例 & 46191 & 64\\
    开发验证 & 四个核心场景，各24例 & 57383 & 96\\
    主测试 & 每个种子下四场景各30例 & 814711、826013、837901 & 360\\
    压力测试 & 五个压力场景，各40例 & 849391 & 200\\
    独立校准 & 四核心场景的独立均匀混合 & 862193 & 240\\
    覆盖评价 & 同分布的独立均匀混合 & 875149 & 320\\
    \bottomrule
  \end{tabular}
\end{table}

四个核心场景分别为静态正常、静态异常、恒速正常和恒速异常。五个压力场景为狭窄几何、加速目标、稀疏观测、标定失配和高污染。主测试采用三个数据种子，以增加独立案例覆盖；这不是三次独立的语言模型研究或规划种子重复，不能据此估计研究控制器的随机性。压力测试有意改变生成机制，用于观察失效边界，不能与同分布主测试合并成一个不加说明的成功率。

冻结记录包含所选配置及其开发验证证据。测试端随后自动产生的比较方法没有反馈给第一轮研究控制器，协议禁止根据其结果修改该冻结方案。早期规则原型的结果已成为第一轮设计的已知背景；本报告后半部分的第二轮又把第一轮结果作为先验，每次后续研究均使用不同的正式测试种子。这里的盲测指工具访问与执行顺序的隔离，并非对拥有完整项目文件的操作者进行密码学保密。可审查性来自冻结文件、协议、输入与预测的保存以及实际事件顺序。

所有数据均由公开模拟器生成。标定失配和高污染只代表预设的合成压力，不是从真实无线设备提取的统计规律。本次没有外部真实数据集或原工业验收数据，任何分场景改善都仅在给定生成机制与样本规模下成立。

\subsection{基线、结构消融与开发前驱}

点估计比较包含四个保留基线：静态WLS、固定参数的静态稳健模型、固定参数的恒速稳健模型，以及第一版的启发式BIC门控。它们由原定位实现包装得到，增加公共诊断不改变已有位置预测。主要比较对象为语言模型最终冻结的方法，因而总体效果同时包含残差、尺度、拟合与模型选择的共同作用。

为区分组件作用，评估脚本依据冻结配置机械生成适用的消融：将对数残差改回距离残差、将常数尺度改回报告尺度、关闭截尾、关闭子集初始化、关闭对数中心修正、将交叉验证门控改为固定恒速，以及将非静态方法改为固定静态。只有所选方案确实包含相应组件时才建立该比较；全部参数变化在协议中保存，且不在测试后重新调优其余参数。

这些比较回答的是“在当前组合附近去掉一个组件会怎样”，不是各个方法族经过公平独立调参后的最优性能比较。例如，更换残差但保持同一稳健尺度，会同时改变该数值尺度的实际含义；若消融表现较差，不能据此断言整个距离残差方法族较差。固定恒速或固定静态也改变计算量，其差异应与工作量共同解释。

脚本还保留登记表中最早两个新增方法作为开发前驱；若配置与最终方法相同则不重复加入。它们不是从最终测试中挑选的最佳中间结果，也不是额外的独立研究运行。前驱比较用于呈现方法修订轨迹，不能替代真实工具日志或单独证明语言模型规划优于其他规划策略。

\subsection{误差统计与配对不确定性}

对每个案例，以当前时刻的位置欧氏误差
\begin{equation}
 e_{m,j}=\|\widehat x_{m,j}(0)-x_j(0)\|_2
\end{equation}
作为基本结果。保存的指标包括均值、RMSE、中位数、P90、P95、最大误差、5米内比例和30米内比例；P90与P95使用NumPy默认的线性分位数插值。文中的CEP90特指这组误差半径的经验90\%分位数。它不等同于原题未定义的“检测精度”，也不能代替静态场景非劣判定。

所有方法使用相同案例。令 $a$ 为冻结方法，$b$ 为比较方法，在场景 $s$ 上定义配对改善
\begin{equation}
 d_{s,j}=e_{b,s,j}-e_{a,s,j},\qquad
 \widehat\Delta=\frac1S\sum_{s=1}^{S}\frac1{n_s}
 \sum_{j=1}^{n_s}d_{s,j}.
 \label{eq:paired-macro-improvement}
\end{equation}
正值表示冻结方法的平均误差较低。\texttt{paired\_ci}在每个核心场景内对差值有放回重采样，再对场景均值等权平均；主测试三个数据种子在各场景内合并，不把三个种子均值当作仅有的三个观测。默认使用种子612303进行3000次重采样，取2.5\%与97.5\%分位数形成名义95\%百分位区间。

这个区间只反映固定方法在当前场景分层与案例抽样下的平均误差差值不确定性。它不是CEP90差异的置信区间，不包含语言模型研究重复、候选选择、实际无线环境迁移或所有计算环境变化。比较多条消融与前驱时，当前实现没有多重比较校正，也没有预注册的单一确证性检验。因此，区间用于探索性的效应量解释；不能把若干区间中偶然较有利的一条称为经过多重检验控制的确认性发现。总体改善仍需结合各静态、移动场景的退化检查。

\subsection{独立校准、覆盖与拒绝曲线}

预测区域仅对已冻结的所选方法进行校准。校准集与覆盖评价集分别独立选择四个核心场景中的一个，再生成案例；这与主测试固定配额的分层设计不同，目的是匹配案例级交换性的条件。标定尺度函数在校准之前已经由定位实现确定，不利用校准真值反复修改。

校准分数与圆盘沿用式\eqref{eq:conformal-score}，考察 $\alpha=0.10$ 和 $\alpha=0.05$ 两个预设水平。实现使用第 $\lceil(m+1)(1-\alpha)\rceil$ 个顺序统计量，而不是线性插值分位数；样本量不足时半径量化因子为正无穷。输出包括总体经验覆盖、平均与中位半径以及分场景覆盖。相应有限样本边际保证依赖交换性与冻结前提\cite{lei2018distributionfree}；一次有限覆盖测试低于名义水平，不会自动否定该边际定理，反过来经验覆盖良好也不能证明前提在真实部署成立。

同一校准因子还应用于五个压力场景，以展示分布变化后的失准，不对该群体宣称分布无关覆盖。拒绝规则的半径门槛预设为2、5、10、15、20、30、50米，分别报告接受率、接受案例数、接受样本平均误差与接受样本覆盖；无接受样本时后两项为空。阈值曲线用于展示覆盖范围与误差的取舍，不作为测试后选择部署阈值的验证依据。尤其不能仅挑一个误差较小但拒绝了大多数案例的点作为总体定位成绩。

当前校准层没有同时为所有基线建立预测圆盘。因此，点估计方法间的不确定性比较由共同案例上的配对区间支持，区域覆盖只评价冻结方法自身；不能据此宣称它的区域比其他方法更窄或更可靠。普通边际覆盖同样不保证拒绝后被接受子集的条件覆盖。

\subsection{计算成本与证据保存}

新求解器的\texttt{work\_units}累计优化器内部实际残差函数调用，包含多个起点、子集求解、截尾轮次、交叉验证折和最终全量重拟合；旧基线保留原求解器记录的函数评价次数。它们都是工作量代理，而非浮点运算总量：一次调用可能处理不同数量的观测，优化器外的评分也不在该计数内。解析Jacobian评价次数与求解次数保存在诊断中，实际耗时另以墙钟时间记录。仅比较方法数量或工具调用次数不足以证明计算预算公平。

基准总耗时包括各方法主测试、压力评价、校准与结果写入，不包含此前的全部语言模型规划、审阅和开发实验成本。科研系统的端到端成本因此需要另读运行时证据，不能把单次定位耗时当作完整研究开销。跨机器耗时差异还受数值库、硬件和系统负载影响。

协议保存冻结文件的字节摘要、定位与工具代码及基准脚本的字节摘要、软件环境、完整输入和逐案例预测。数值求解若发生无法处理的错误，当前评估会中止而非悄悄删除案例；只有全部指定案例完成后，相关汇总才构成完整评估。审计关注的是配置、数据、输出和结论能否相互对应，而不是要求时间戳或不同硬件上的耗时逐字节一致。

\section{第一轮真实智能体运行与研究迭代}\label{sec:trace}

\subsection{运行身份、预算与实际调用}

正式轨迹保存在\texttt{runs/live\_reference/}。本次状态为\texttt{frozen}，共执行13个研究动作，发起14次模型调用，其中14次成功，包含1次独立审阅调用。实验工具调用计数为6，其中控制消融2次。正式过程历时31.6分钟。

模型由已登录的Codex CLI调用，采用用户配置的默认模型，没有人为填写服务端未返回的型号。本次保存的响应不是回放。服务内部连接重试与控制器重新请求是不同计数；完整日志均保留。

\begin{table}[H]\centering\small
\caption{正式研究的审计计数。token使用量按CLI返回字段相加，不推断费用。}

\begin{tabular}{lr}
\toprule
指标 & 数值\\
\midrule
注册的新方法（含消融） & 5\\
成功模型调用 & 14\\
独立审阅调用 & 1\\
输入token & 444020\\
缓存输入token & 30720\\
输出token & 6994\\
检测到的CLI额外工具事件 & 0\\
\bottomrule\end{tabular}\end{table}

\subsection{由一次失败驱动的真实修订}

模型最初预期对数残差有助于移动异常案例，但P2的开发均值达到12.117米，并未实现预期。随后它调用诊断工具检查极端误差。工具返回的是跨场景全局最差案例，因此模型明确承认，原先关于某个场景最差四例的假设尚不能完整判定；它没有将不匹配的诊断输出当成已证实的因果解释。

第8步只将拟合损失从soft-L1改为Cauchy，登记P4；观测划分、残差、尺度和验证评分保持不变。工具返回的移动异常均值从12.117降至8.555米，CEP90从24.954降至14.923米。然而静态正常均值为1.384米，高于事前开发判据1.316米，故这个包含静态条件的联合假设仍被否定。10\%容忍度是模型在开发阶段声明的经验判据，不是原工业要求的严格非劣，也不是统计非劣检验。

这一链条中的假设、诊断对象、修改与拒绝理由来自实际模型响应；上述数值来自保存的配对实验。因而可以核查模型是否根据反馈改变研究动作，而不是仅阅读最后一个最优配置。

\noindent\textbf{证据定位：}研究目录下的\path{tools/experiments/call_003_discovery_log_cv_reported.json}、\path{tools/experiments/call_003_discovery_log_cv_cauchy.json}，以及对应步号的\path{events.jsonl}。

\subsection{逐步动作与可核查证据}

下表由事件日志自动生成。方法代号按首次注册顺序分配；它们不是预设候选排名。实际计算结果来自工具返回，研究动机与预期来自模型响应。

\begingroup\small\setlength{\tabcolsep}{4pt}

\begin{longtable}{@{}p{8mm}p{23mm}p{121mm}@{}}

\caption{实际研究动作序列。均值单位为米；开发与验证来自不同案例集。}\label{tab:trace}\\

\toprule 步 & 动作 & 执行内容与实际结果\\\midrule\endfirsthead

\toprule 步 & 动作 & 执行内容与实际结果\\\midrule\endhead

\bottomrule\endfoot

1 & 读取来源 & 目前尚无来源证据，先读取问题材料、假设与算子的精确定义，为基线实验和可检验方法设计建立依据。\\[3pt]

2 & 执行实验 & 开发集：静态WLS 均值 6.599；静态稳健 均值 5.868；恒速稳健 均值 4.060；旧BIC门控 均值 3.896\\[3pt]

3 & 注册方法 & P1；距离域soft-L1观测交叉验证选择在discovery上同时满足：moving\_\allowbreak{}outliers平均误差低于legacy\_\allowbreak{}bic，static\_\allowbreak{}clean平均误差不超过static\_\allowbreak{}wls的1.10倍。\\[3pt]

4 & 注册方法 & P2；在discovery共同案例上，log\_\allowbreak{}cv\_\allowbreak{}reported的moving\_\allowbreak{}outliers平均当前位置误差低于range\_\allowbreak{}cv\_\allowbreak{}reported，且static\_\allowbreak{}clean平均误差不超过static\_\allowbreak{}wls的1.10倍。\\[3pt]

5 & 注册方法 & P3；在discovery共同案例上，range\_\allowbreak{}cv\_\allowbreak{}constant的moving\_\allowbreak{}outliers平均当前位置误差低于range\_\allowbreak{}cv\_\allowbreak{}reported，且static\_\allowbreak{}clean平均误差不超过static\_\allowbreak{}wls的1.10倍。\\[3pt]

6 & 执行实验 & 开发集：静态WLS 均值 6.599；P1 均值 3.943；P2 均值 5.209；P3 均值 5.329\\[3pt]

7 & 失败诊断 & 诊断 P2；discovery支持距离域交叉验证候选的预设联合条件，但对数残差和共享尺度的移动异常改善假设均被否定；共享尺度还违反正常静止容忍度。先诊断对数候选的严重失败，区分模型误选与拟合异常，为修订提供证据。\\[3pt]

8 & 控制消融 & 开发集：P2 均值 5.209；P4 均值 4.294；改变 \{'loss': 'cauchy'\}\\[3pt]

9 & 数学检查 & Cauchy消融降低了移动异常均值和尾部误差，但正常静止均值1.38382米超过预设1.31631米阈值，因此联合假设被否定。下一步检查数学命题，明确观测权重偏差、几何不可辨识性及不确定性校准的适用条件。\\[3pt]

10 & 控制消融 & 开发集：P1 均值 3.943；P5 均值 4.058；改变 \{'model': 'velocity'\}\\[3pt]

11 & 执行实验 & 验证集：静态WLS 均值 6.022；旧BIC门控 均值 3.912；P1 均值 3.742；P5 均值 4.759\\[3pt]

12 & 执行实验 & 验证集：静态WLS 均值 6.022；静态稳健 均值 5.572；恒速稳健 均值 4.759；P1 均值 3.742\\[3pt]

13 & 冻结方法 & 冻结 P1；冻结已登记且未据validation修改的range\_\allowbreak{}cv\_\allowbreak{}reported。验证每场景24例：四场景等权平均误差3.74151米，低于static\_\allowbreak{}robust的5.57182、dynamic\_\allowbreak{}robust的4.75875及legacy\_\allowbreak{}bic的3.91248米；正常静止0.90199米低于预设WLS均值1.10倍阈值1.01550米；移动异常8.36549米低于BIC的9.\\[3pt]

\end{longtable}\endgroup

\subsection{候选组成、修订与审阅}

\begin{table}[H]\centering\small
\caption{智能体注册的方法组成与已观测均值。空白项表示未调用该划分。}

\begin{tabular}{p{8mm}p{104mm}rr}
\toprule
代号 & 模型／残差／尺度／损失 & 开发 & 验证\\
\midrule
P1 & \texttt{cv\_\allowbreak{}select/range/reported/soft\_\allowbreak{}l1} \newline trim=0.0, subsets=0, margin=0.05 & 3.943 & 3.742\\
P2 & \texttt{cv\_\allowbreak{}select/log\_\allowbreak{}range/reported/soft\_\allowbreak{}l1} \newline trim=0.0, subsets=0, margin=0.05 & 5.209 & --\\
P3 & \texttt{cv\_\allowbreak{}select/range/constant/soft\_\allowbreak{}l1} \newline trim=0.0, subsets=0, margin=0.05 & 5.329 & --\\
P4 & \texttt{cv\_\allowbreak{}select/log\_\allowbreak{}range/reported/cauchy} \newline trim=0.0, subsets=0, margin=0.05 & 4.294 & --\\
P5 & \texttt{velocity/range/reported/soft\_\allowbreak{}l1} \newline trim=0.0, subsets=0, margin=0.05 & 4.058 & 4.759\\
\bottomrule\end{tabular}\end{table}

\paragraph{审阅1（第3次普通实验后）}

方法登记存在定义冲突：log\_\allowbreak{}cv\_\allowbreak{}cauchy仍写soft-L1目标，range\_\allowbreak{}velocity\_\allowbreak{}reported仍写交叉验证选择。须明确实际执行的spec与数学定义，避免错误归因。；验证联合判据按点估计全部通过：移动异常8.36549<9.22411米，正常静止0.90199<1.01550米，宏均值3.74151<4.75930米；但每场景仅24例，缺少配对差异的不确定性，不能称为统计非劣或泛化保证。；计算预算尚未超限，但比较未匹配计算量：候选平均69.35工作单位，BIC为47.49，强制恒速为26.49。调用次数不能替代实际计算预算；应报告精度与成本权衡。；验证缺少static\_\allowbreak{}robust和dynamic\_\allowbreak{}robust；尤其匹配配置的强制静态稳健基线有助于分离稳健损失与模型选择的贡献。；未见测试或校准泄漏证据，但来源、划分独立性及预处理实现未提供，无法独立核实。交错观测验证评价的是留出距离预测，与当前位置估计目标并不等价。；多候选与诊断修订均依据discovery，属于适当开发；继续根据已见validation调参会产生选择偏差。宏均值结论还依赖四场景等权设定。；最差4例诊断实际跨场景，尚未完成原先针对移动异常最差4例的假设检验。镜像反例否定局部满秩足以保证全局可辨识；恒速基线在移动异常上更优，也否定候选普遍占优。

建议的下一步：先澄清两项登记定义冲突并记录执行配置；保持候选不变，用一次集中validation调用补齐static\_\allowbreak{}robust、dynamic\_\allowbreak{}robust及候选比较。随后依据预设静止容忍度、场景权重及计算成本冻结有验证证据的方法，明确未解决校准问题；不要为达到建议步数继续调参。

\paragraph{对审阅意见的实际回应}

模型在第12步的持久记忆中记录：方法ID不可修改，保留并明确登记文字缺陷：log\_\allowbreak{}cv\_\allowbreak{}cauchy的spec为cauchy，拟合损失应为c²log(1+(r/c)²)，留出评分仍为单位尺度soft-L1；range\_\allowbreak{}velocity\_\allowbreak{}reported的spec为velocity，不执行CV选择。以上依据工具返回的spec及受控变更，未独立核验实现。主要候选不变，本次补齐稳健基线，不据validation调参；随后依据既定静止10\%容忍度、四场景等权误差及成本选择冻结。现有点估计不是统计非劣证明，基线配置未确认一致前不能视为单因素消融；覆盖率、配对不确定性及工业适用性仍未解决。

因此，本报告的方法数学解释以实际执行配置和数值源码为准。历史登记文本被原样保留，不能把上述文字纠错说成重新运行了不同的定位算法；后续版本应将父方法说明与消融后的定义分开保存。

\subsection{最终冻结配置}

最终方法为P1，原始标识为\texttt{range\_\allowbreak{}cv\_\allowbreak{}reported}。参数和结论在最终测试前保存；下列配置不因测试结果而更改。

\begin{table}[H]\centering\small
\caption{冻结方法的实际声明参数；未声明字段采用数值模块公开默认值。}

\begin{tabular}{ll}
\toprule
字段 & 值\\
\midrule
\texttt{consensus\_\allowbreak{}trials} & \texttt{0}\\
\texttt{f\_\allowbreak{}scale} & \texttt{1.5}\\
\texttt{gate\_\allowbreak{}margin} & \texttt{0.05}\\
\texttt{log\_\allowbreak{}bias\_\allowbreak{}correction} & \texttt{True}\\
\texttt{log\_\allowbreak{}sd} & \texttt{0.1}\\
\texttt{loss} & \texttt{soft\_\allowbreak{}l1}\\
\texttt{model} & \texttt{cv\_\allowbreak{}select}\\
\texttt{residual} & \texttt{range}\\
\texttt{ridge} & \texttt{0.05}\\
\texttt{scale} & \texttt{reported}\\
\texttt{starts} & \texttt{3}\\
\texttt{trim\_\allowbreak{}fraction} & \texttt{0.0}\\
\bottomrule\end{tabular}\end{table}

冻结时模型给出的研究结论为：冻结已登记且未据validation修改的range\_\allowbreak{}cv\_\allowbreak{}reported。验证每场景24例：四场景等权平均误差3.74151米，低于static\_\allowbreak{}robust的5.57182、dynamic\_\allowbreak{}robust的4.75875及legacy\_\allowbreak{}bic的3.91248米；正常静止0.90199米低于预设WLS均值1.10倍阈值1.01550米；移动异常8.36549米低于BIC的9.22411米，但高于dynamic\_\allowbreak{}robust的6.08423米。候选平均69.35417工作单位，高于静态稳健20.53125、动态稳健26.95833及BIC的47.48958；比较未匹配计算量，选择体现预设静止容忍度与四场景等权误差的折中，不代表普遍占优。discovery拒绝了对数残差和常数尺度分别优于距离域reported的假设；Cauchy修订改善对数方法移动异常误差但未通过静止容忍度。受控强制恒速消融显示移动收益与静止代价。登记文字缺陷须保留说明：log\_\allowbreak{}cv\_\allowbreak{}cauchy实际spec为cauchy，其拟合损失为c²log(1+(r/c)²)，留出评分仍为单位尺度soft-L1；range\_\allowbreak{}velocity\_\allowbreak{}reported实际spec为velocity，不执行CV选择；这些解释依工具spec，未经独立实现核验。数学检查支持指定标量乘性噪声模型的自加权偏差公式及镜像不可辨识恒等式，不可推广为完整定位误差证明；局部满秩不足以排除全局歧义。共形边际覆盖需要估计器及正尺度冻结和校准测试交换性，尚无校准或覆盖率实测。仍缺配对差异不确定性、独立数据与实现核验，原移动异常最差4例诊断假设未完成；观测预测评分与当前位置误差目标不等价。所有结论限于已见合成开发数据的点估计，不声称统计非劣、工业验收、原创性或泛化保证。冻结后立即结束开发，测试及校准尚未访问。
\section{第一轮独立评测、消融与失败分析}\label{sec:results}

\subsection{常规测试的总体误差与对照}

本次基准累计完成5040次定位预测，壁钟时间为0.8分钟。主测试360例采用三个独立数据种子，各常规场景样本数相同。所有方法使用同一案例；所选配置在评测前已冻结。

\begin{table}[H]\centering\small
\caption{独立主测试总体性能。CEP90是案例位置误差的经验90\%分位数。}

\begin{tabular}{lrrrrr}
\toprule
方法 & 均值/m & 中位/m & CEP90/m & RMSE/m & 平均工作量\\
\midrule
静态WLS & 6.532 & 5.340 & 14.959 & 8.504 & 19.3\\
静态稳健 & 6.196 & 3.928 & 14.993 & 8.395 & 20.6\\
恒速稳健 & 4.674 & 2.993 & 9.853 & 6.781 & 27.2\\
旧BIC门控 & 4.347 & 2.142 & 12.203 & 6.790 & 47.8\\
智能体所选 & 4.189 & 2.252 & 9.966 & 6.434 & 71.3\\
\bottomrule\end{tabular}\end{table}

\begin{table}[H]\centering\small
\caption{场景内配对bootstrap比较：对照误差减去所选方法误差。正值表示所选方法较低。}

\begin{tabular}{lrr}
\toprule
对照 & 均值减少/m & 95\%区间/m\\
\midrule
静态WLS & 2.343 & [1.890, 2.776]\\
静态稳健 & 2.007 & [1.578, 2.424]\\
恒速稳健 & 0.485 & [0.064, 0.912]\\
旧BIC门控 & 0.158 & [-0.091, 0.425]\\
\bottomrule\end{tabular}\end{table}

区间描述固定生成机制下的有限样本差异，没有对众多方法比较作多重校正。它不能度量重新运行一次大模型规划器时的方差，也不能据此归因规划器优于其它研究策略。

\subsection{逐场景结果与误差分布}

\begin{table}[H]\centering\small
\caption{四类常规场景的平均当前位置误差（m），每类90例。}

\begin{tabular}{lrrrr}
\toprule
方法 & 静态正常 & 静态异常 & 移动正常 & 移动异常\\
\midrule
静态WLS & 1.031 & 3.492 & 10.349 & 11.259\\
静态稳健 & 1.059 & 2.080 & 10.429 & 11.217\\
恒速稳健 & 2.636 & 6.229 & 2.619 & 7.213\\
旧BIC门控 & 1.305 & 2.080 & 2.984 & 11.019\\
智能体所选 & 1.419 & 2.329 & 3.681 & 9.328\\
\bottomrule\end{tabular}\end{table}

\begin{figure}[H]\centering
\begin{tikzpicture}\begin{axis}[ybar,bar width=9pt,width=0.94\linewidth,height=6.3cm,
ylabel={平均位置误差（m）},ymin=0,grid=major,enlarge x limits=0.15,
xtick={0,1,2,3},xticklabels={静态正常,静态异常,移动正常,移动异常},
legend style={at={(0.5,1.04)},anchor=south,legend columns=3},
x tick label style={font=\small,text width=17mm,align=center},]
\addplot[fill=navy,draw=navy] coordinates {(0,1.030812) (1,3.491518) (2,10.348803) (3,11.258691)};
\addlegendentry{静态WLS}
\addplot[fill=orange,draw=orange] coordinates {(0,2.635838) (1,6.229444) (2,2.619052) (3,7.212711)};
\addlegendentry{恒速稳健}
\addplot[fill=teal,draw=teal] coordinates {(0,1.418681) (1,2.329268) (2,3.681274) (3,9.327780)};
\addlegendentry{智能体所选}
\end{axis}\end{tikzpicture}
\caption{常规场景的配对误差比较。总体改善不能替代逐场景非劣分析。}
\label{fig:main-bars}
\end{figure}

\begin{figure}[H]\centering
\begin{tikzpicture}\begin{axis}[width=0.92\linewidth,height=7cm,xlabel={位置误差（m）},ylabel={经验累计概率},xmin=0,xmax=30,ymin=0,ymax=1,grid=major,legend pos=south east]\addplot[thick,no marks,color=navy] table[x=error,y=cdf] {data/cdf_static_wls.dat};\addlegendentry{静态WLS}
\addplot[thick,no marks,color=orange] table[x=error,y=cdf] {data/cdf_static_robust.dat};\addlegendentry{静态稳健}
\addplot[thick,no marks,color=brick] table[x=error,y=cdf] {data/cdf_dynamic_robust.dat};\addlegendentry{恒速稳健}
\addplot[thick,no marks,color=gray] table[x=error,y=cdf] {data/cdf_legacy_bic.dat};\addlegendentry{旧BIC门控}
\addplot[thick,no marks,color=teal] table[x=error,y=cdf] {data/cdf_llm_selected.dat};\addlegendentry{智能体所选}\end{axis}\end{tikzpicture}
\caption{独立主测试误差的经验分布。横轴显示0--30米；超过30米的样本仍保留在全部统计中。}
\label{fig:cdf}
\end{figure}

\subsection{组件消融与开发前驱}

\begin{table}[H]\centering\small
\caption{冻结方法的控制消融与开发前驱。消融不重新调参。}

\begin{tabular}{lrrrr}
\toprule
方法 & 主测试均值 & 主测试CEP90 & 压力均值 & 与所选均值差\\
\midrule
固定恒速模型 & 4.674 & 9.850 & 8.814 & 0.485\\
固定静态模型 & 6.196 & 14.994 & 12.839 & 2.007\\
开发前驱1 & 4.704 & 12.564 & 12.396 & 0.515\\
\bottomrule\end{tabular}\end{table}

某个消融若优于冻结方法，应作为预先冻结方案的负面证据报告，不能用该测试结果悄悄替换最终方法。未触发的消融表示冻结配置本来没有该组件，避免重复计算伪造独立对照。

\subsection{分布外压力测试}

每类压力场景独立生成40例，改变几何、运动、观测数量、标定或污染条件。下表同时保留均值与尾部误差，以检验常规场景的改善能否延续到这些变化。

\begin{table}[H]\centering\small
\caption{压力场景的均值／CEP90（m），每类40例。}

\begin{tabular}{lrrr}
\toprule
场景 & 静态WLS & 恒速稳健 & 智能体所选\\
\midrule
狭窄几何 & 14.503 / 21.270 & 9.817 / 24.461 & 9.432 / 22.286\\
加速运动 & 10.846 / 15.653 & 4.414 / 7.723 & 5.313 / 11.554\\
稀疏观测 & 10.805 / 17.493 & 5.976 / 14.303 & 11.121 / 18.090\\
标定漂移 & 15.044 / 21.988 & 11.221 / 21.014 & 13.278 / 24.098\\
高比例污染 & 10.862 / 21.049 & 12.641 / 25.044 & 11.734 / 22.187\\
\bottomrule\end{tabular}\end{table}

\begin{figure}[H]\centering
\begin{tikzpicture}\begin{axis}[ybar,bar width=9pt,width=0.94\linewidth,height=6.3cm,
ylabel={平均位置误差（m）},ymin=0,grid=major,enlarge x limits=0.15,
xtick={0,1,2,3,4},xticklabels={狭窄几何,加速运动,稀疏观测,标定漂移,高比例污染},
legend style={at={(0.5,1.04)},anchor=south,legend columns=3},
x tick label style={font=\small,text width=17mm,align=center},]
\addplot[fill=navy,draw=navy] coordinates {(0,14.502767) (1,10.845790) (2,10.804917) (3,15.044343) (4,10.861705)};
\addlegendentry{静态WLS}
\addplot[fill=orange,draw=orange] coordinates {(0,9.817206) (1,4.414207) (2,5.976103) (3,11.221198) (4,12.640590)};
\addlegendentry{恒速稳健}
\addplot[fill=teal,draw=teal] coordinates {(0,9.431836) (1,5.312884) (2,11.121400) (3,13.277567) (4,11.734202)};
\addlegendentry{智能体所选}
\end{axis}\end{tikzpicture}
\caption{冻结估计器在五类压力条件下的误差。压力输入不参与方法选择。}
\label{fig:stress-bars}
\end{figure}

\begin{table}[H]\centering\small
\caption{所选方法实际使用的运动模型次数。标签来自求解器诊断。}

\begin{tabular}{lrrr}
\toprule
场景 & 静态 & 恒速 & 其它\\
\midrule
静态正常 & 78 & 12 & 0\\
静态异常 & 85 & 5 & 0\\
移动正常 & 12 & 78 & 0\\
移动异常 & 49 & 41 & 0\\
狭窄几何 & 11 & 29 & 0\\
加速运动 & 5 & 35 & 0\\
稀疏观测 & 40 & 0 & 0\\
标定漂移 & 21 & 19 & 0\\
高比例污染 & 30 & 10 & 0\\
\bottomrule\end{tabular}\end{table}

\subsection{独立校准、覆盖与拒绝}

使用240个独立混合场景案例校准尺度，再在320个同分布案例及200个压力案例上检查覆盖。名义水平来自预定协议，实测覆盖与区域大小均由已保存的逐例结果计算。

\begin{table}[H]\centering\small
\caption{独立校准后的实际覆盖。压力分布的覆盖没有交换性保证。}

\begin{tabular}{lrrrrrr}
\toprule
名义覆盖 & 分位数$q$ & 评估总体 & 样本 & 实际覆盖 & 均值半径/m & 中位半径/m\\
\midrule
90.0\% & 3.548 & 同分布混合 & 320 & 85.6\% & 8.481 & 7.084\\
90.0\% & 3.548 & 压力分布 & 200 & 74.0\% & 20.147 & 14.300\\
95.0\% & 5.199 & 同分布混合 & 320 & 94.7\% & 12.429 & 10.383\\
95.0\% & 5.199 & 压力分布 & 200 & 85.5\% & 29.527 & 20.958\\
\bottomrule\end{tabular}\end{table}

\begin{table}[H]\centering\small
\caption{90\%名义区域的拒绝曲线：只接受半径不超过阈值的案例。}

\begin{tabular}{lrrrr}
\toprule
半径阈值/m & 接受数 & 接受率 & 接受均值误差/m & 接受后覆盖\\
\midrule
2 & 0 & 0.0\% & -- & --\\
5 & 85 & 26.6\% & 2.034 & 88.2\%\\
10 & 229 & 71.6\% & 3.165 & 83.8\%\\
15 & 291 & 90.9\% & 3.540 & 84.5\%\\
20 & 311 & 97.2\% & 3.702 & 85.2\%\\
30 & 318 & 99.4\% & 3.870 & 85.5\%\\
50 & 320 & 100.0\% & 3.955 & 85.6\%\\
\bottomrule\end{tabular}\end{table}

拒绝阈值越小通常会减少输出数量。表中的接受后覆盖是实际样本统计，不继承总体边际覆盖定理；缺失的接受子集以空值表示。较大的预测圆盘即使覆盖充分，也可能缺乏实际定位价值，因此必须同时报告半径与接受率。

\subsection{独立数学见证}

在标量检验中，真距离为20米，对数噪声标准差为0.35。解析自加权最优值为15.65409米，200000次独立模拟得到15.67145米；修正后对数目标给出20.01864米。该见证用于核对解析公式，不构成新的泛化实验。

镜像构造的两条轨迹当前位置相差6.0米，所有距离差的最大绝对值为0.0e+00米；数据Jacobian秩为4，最小奇异值为0.96509。因此，这个实例直接展示了局部满秩与全局镜像歧义能够共存。

\subsection{哪些结论成立，哪些预期没有泛化}

\paragraph{总体均值有改善，不能据此宣称处处更优。}
冻结的距离域交叉验证方法在360例主测试上的平均误差为4.189米，相比WLS的6.532米降低35.87\%。它比固定恒速稳健基线的均值低10.38\%，但CEP90为9.966米，略高于后者的9.853米。相对旧BIC门控，均值只减少0.158米，配对95\%区间为$[-0.091,0.425]$米，包含零；本轮不能确认其稳定优于BIC。一个总体均值无法同时概括静态、动态与尾部风险。

\paragraph{静止非劣条件没有从验证集推广到测试集。}
验证集静态正常误差为0.902米，低于同批WLS的0.923米；但独立测试中变为1.419米，WLS为1.031米，退化37.63\%。因此，即使把开发判据放宽为10\%容忍度，测试结果仍未满足。更不能把总体35.87\%的均值改善写成完成原题的“标准场景不劣化”或“移动检测精度提高20\%”。原题的后一指标没有得到定义，也没有真实工业对照数据。

\paragraph{门控节省了静态过拟合的代价，也丢失了一部分运动收益。}
固定恒速消融的主测试均值为4.674米，高于所选方案，但其压力测试均值为8.814米，低于所选方案的10.176米。所选方案在稀疏场景的平均误差为11.121米，固定恒速基线为5.976米。这里有可核查的实现机制：三折划分下，8份观测产生不足6条的训练折，运动模型因预设可用性门槛被排除。它不是从数据中“发现目标静止”；这种保守回退需要作为系统失效方式报告。

\paragraph{计算成本提高，稳健性仍有明确边界。}
主测试上所选方法平均消耗71.3个工作单位，旧BIC为47.8、恒速稳健为27.2；本机记录的平均用时约为14.9、7.6和4.9毫秒。工作单位在新旧求解器之间的计数口径有差异，壁钟时间也依赖实现和设备，所以这些数字只说明本次运行的精度与成本取舍。最终配置没有使用截尾或子集初始化，不能把本轮效果归因于这些虽然已实现、但没有进入冻结方案的组件。

\paragraph{不确定性层必须报告实测欠覆盖。}
名义90\%区域在独立同分布混合样本上仅覆盖274/320例，即85.625\%；名义95\%区域覆盖303/320例，即94.6875\%。移动异常子群的相应覆盖为59.76\%和87.80\%，静态两个子群则均为100\%。这展示了总体校准与子群均衡之间的差别。压力分布上的覆盖进一步降至74.0\%和85.5\%，尤其稀疏场景分别只有37.5\%和60.0\%。

共形定理对随机校准样本与新案例联合取概率，不承诺任一已实现校准集下的有限测试比例必定超过名义水平。本次90\%经验覆盖没有达标，这一事实保留；既不能将它掩盖为“90\%覆盖已获验证”，也不能仅凭一个有限样本比例否定带有条件的边际覆盖定理。半径10米的90\%区域拒绝策略在同分布测试中接受71.56\%案例，接受后的覆盖只有83.84\%，再次说明按半径拒绝不会自动保留条件覆盖。

上述结果支持一个有限结论：本研究得到一个对常规混合场景有平均收益的可审计方案，并定位了静止泛化、稀疏门控、尾部风险与子群校准方面的失败。开放问题仍然存在，失败结果本身为下一轮建模提供了更明确的约束。


\subsection{校准随机性的事后数学参考}
\label{sec:coverage-variability-reference}

名义90\%区域在本次独立覆盖检验中的经验覆盖率为
$274/320=85.625\%$，这一不足必须保留。以下计算只解释有限校准样本的随机性，
不是新增实验、事后重新校准或方法修订，也不把未达名义水平改写成成功。

假设位置估计器与正尺度函数均已在校准前冻结，校准分数和测试分数独立同分布，
且其共同累积分布函数$F$连续。令校准样本数$m=240$，
$q=S_{(k)}$，$k=\lceil(m+1)(1-\alpha)\rceil$。
概率积分变换与次序统计量给出
\[
 U=F(q)\sim\operatorname{Beta}(k,m+1-k),\qquad
 K\mid U\sim\operatorname{Binomial}(320,U),
\]
其中$K$为320个新案例中被区域覆盖的数量。对随机$U$积分后，
得到
\[
 K\sim\operatorname{BetaBinomial}(320,k,m+1-k).
\]
这是\emph{校准集和测试集共同重新抽样}的无条件分布；
固定本次校准集后的覆盖概率$F(q)$仍未知，因此下表不是该条件覆盖率的置信区间。

\begin{table}[H]
\centering\small
\caption{理想连续IID评分模型下的事后覆盖计数参考。预测区间取离散分布的2.5\%和97.5\%分位数。}
\label{tab:coverage-variability-reference}
\begin{tabular}{@{}crrrcr@{}}
\toprule
名义覆盖 & $k$ & 实测$K$ & $\mathbb{E}[K]$ & 95\%预测区间 & $\Pr(K\le K_{\rm obs})$\\
\midrule
90\% & 217 & 274 & 288.13 & $[271,303]$ & 0.0563\\
95\% & 229 & 303 & 304.07 & $[291,314]$ & 0.4258\\
\bottomrule
\end{tabular}
\end{table}

对应的覆盖比例预测区间分别为$[84.6875\%,94.6875\%]$与
$[90.9375\%,98.125\%]$。数值使用\texttt{scipy.stats.betabinom}计算，
参数、软件版本及完整小数保存在
\nolinkurl{research/coverage_variability_reference.json}。
这些参考说明本次较低覆盖与理想联合抽样模型并不必然矛盾，但不能验证实际
评分分布的连续IID假设、保证给定校准集的覆盖，或弥补已观察到的85.625\%覆盖率。
它们也不提供分场景、拒绝子集或分布偏移下的保证；研究未据此调整$q$、
尺度函数、阈值或定位方法。

\section{第一轮的观测数与污染敏感性研究}\label{sec:sensitivity}

\subsection{观测数量与污染水平的预登记敏感性分析}
\label{sec:sensitivity-protocol}

为补充单一压力点的评价，另以独立脚本\texttt{sensitivity.py}研究适用范围。
在生成数据前保存的\texttt{research/sensitivity\_plan.json}固定根种子980221，
观测数$n\in\{8,12,16,24\}$、名义污染比例
$\epsilon\in\{0,0.1,0.25,0.4\}$，以及静态异常、恒速异常两个场景。
每个网格单元独立生成12例，合计32格、384例。各格使用由场景编号、观测数和
污染水平派生的独立伪随机子流；三种方法在每格使用完全相同的案例，因此对照差异
按案例配对。固定比较方法仅为最终冻结方案、恒速稳健基线和静态加权最小二乘，
总计1152次预测，不依此重新选择方法或修改参数。

脚本先验证研究冻结、关键源码与开发证据，再生成数据。协议保存预登记文件、
本脚本及冻结方法的摘要；这些本地记录用于审计顺序，不等于第三方时间戳认证。
每格报告平均位置误差、经验CEP90、计算工作量、壁钟时间和实际模型选择频率，
同时保存逐例输入、预测与CV诊断。污染数采用生成器的
$\operatorname{round}(n\epsilon)$，因而小样本下的实际比例具有离散性；
结果必须并列记录名义比例与实际比例。

观测数同时影响信息量与CV门控。默认三折时间交错划分在$n=8$时含训练样本数
为5的折，低于运动模型要求的6份训练观测；若冻结方案采用该CV配置，运动模型会
因可用性条件被排除。该机制可由完成折数与运动可用标记核对，不能将其误读为数据
支持目标静止。若最终方案不含CV，则相关字段记为不适用，不补造门控结果。
本分析仅描述预设合成网格，每格12例不足以精确定位稳健性相变；不作逐格显著性
或多重比较确证，不重新校准误差区域，也不提出新的覆盖率保证。


实际完成384个独立案例、1152次定位预测，耗时0.2分钟。以下结果没有返馈给第一轮研究智能体，也没有用于替换其已冻结方案；第二轮则明确把这些结果作为已暴露先验，另用新数据评价修订。

在全部32个预设网格单元中，所选方法的样本均值低于静态WLS的单元为23个，低于恒速稳健基线的单元为20个。每格仅12例，单元胜出数量是描述性结果，不是显著性检验或部署分布下的成功概率。

\begin{table}[H]\centering\small
\caption{静止目标：所选方法平均位置误差（m）。列为名义污染比例。}

\begin{tabular}{lrrrr}
\toprule
观测数 & 0\% & 10\% & 25\% & 40\%\\
\midrule
8 & 1.807 & 3.042 & 2.810 & 4.937\\
12 & 2.295 & 2.388 & 2.261 & 11.469\\
16 & 1.189 & 2.861 & 2.306 & 3.184\\
24 & 0.897 & 1.506 & 1.940 & 6.133\\
\bottomrule\end{tabular}\end{table}

\begin{figure}[H]\centering
\begin{tikzpicture}\begin{axis}[width=0.91\linewidth,height=6.4cm,xlabel={观测数},ylabel={平均位置误差（m）},xtick={8,12,16,24},ymin=0,grid=major,legend style={at={(0.5,1.04)},anchor=south,legend columns=4}]
\addplot[thick,color=navy,mark=*] coordinates {(8,1.806937) (12,2.295357) (16,1.189295) (24,0.896926)};\addlegendentry{0\%}
\addplot[thick,color=teal,mark=square*] coordinates {(8,3.041615) (12,2.387690) (16,2.861054) (24,1.506256)};\addlegendentry{10\%}
\addplot[thick,color=orange,mark=triangle*] coordinates {(8,2.809855) (12,2.261038) (16,2.305569) (24,1.940370)};\addlegendentry{25\%}
\addplot[thick,color=brick,mark=diamond*] coordinates {(8,4.936753) (12,11.468595) (16,3.184007) (24,6.133056)};\addlegendentry{40\%}
\end{axis}\end{tikzpicture}
\caption{静止目标的样本量与污染敏感性。线段仅连接离散实验格点，不表示格点间的理论插值。}
\label{fig:sensitivity-static_outliers}
\end{figure}

\begin{table}[H]\centering\small
\caption{静止目标：模型选择及相对恒速基线的配对均值差。正差值表示所选方法误差较低。}

\begin{tabular}{lrrrrr}
\toprule
观测数 & 名义污染 & 实际污染 & 选择恒速 & CV可用 & 误差减少/m\\
\midrule
8 & 0\% & 0.0\% & 0.0\% & 0.0\% & 3.582\\
8 & 10\% & 12.5\% & 0.0\% & 0.0\% & 11.203\\
8 & 25\% & 25.0\% & 0.0\% & 0.0\% & 11.262\\
8 & 40\% & 37.5\% & 0.0\% & 0.0\% & 16.244\\
12 & 0\% & 0.0\% & 25.0\% & 100.0\% & 1.646\\
12 & 10\% & 8.3\% & 25.0\% & 100.0\% & 1.097\\
12 & 25\% & 25.0\% & 0.0\% & 100.0\% & 4.197\\
12 & 40\% & 41.7\% & 41.7\% & 100.0\% & 6.950\\
16 & 0\% & 0.0\% & 0.0\% & 100.0\% & 1.332\\
16 & 10\% & 12.5\% & 16.7\% & 100.0\% & 2.038\\
16 & 25\% & 25.0\% & 0.0\% & 100.0\% & 9.356\\
16 & 40\% & 37.5\% & 8.3\% & 100.0\% & 7.023\\
24 & 0\% & 0.0\% & 0.0\% & 100.0\% & 0.964\\
24 & 10\% & 8.3\% & 16.7\% & 100.0\% & 1.237\\
24 & 25\% & 25.0\% & 0.0\% & 100.0\% & 1.827\\
24 & 40\% & 41.7\% & 16.7\% & 100.0\% & 5.043\\
\bottomrule\end{tabular}\end{table}

\begin{table}[H]\centering\small
\caption{移动目标：所选方法平均位置误差（m）。列为名义污染比例。}

\begin{tabular}{lrrrr}
\toprule
观测数 & 0\% & 10\% & 25\% & 40\%\\
\midrule
8 & 11.526 & 9.125 & 10.956 & 14.169\\
12 & 4.647 & 6.090 & 10.801 & 14.661\\
16 & 4.661 & 12.936 & 10.409 & 9.317\\
24 & 1.778 & 2.938 & 5.347 & 8.704\\
\bottomrule\end{tabular}\end{table}

\begin{figure}[H]\centering
\begin{tikzpicture}\begin{axis}[width=0.91\linewidth,height=6.4cm,xlabel={观测数},ylabel={平均位置误差（m）},xtick={8,12,16,24},ymin=0,grid=major,legend style={at={(0.5,1.04)},anchor=south,legend columns=4}]
\addplot[thick,color=navy,mark=*] coordinates {(8,11.526413) (12,4.647328) (16,4.661318) (24,1.777933)};\addlegendentry{0\%}
\addplot[thick,color=teal,mark=square*] coordinates {(8,9.125434) (12,6.090254) (16,12.936288) (24,2.938467)};\addlegendentry{10\%}
\addplot[thick,color=orange,mark=triangle*] coordinates {(8,10.955595) (12,10.801396) (16,10.409263) (24,5.347258)};\addlegendentry{25\%}
\addplot[thick,color=brick,mark=diamond*] coordinates {(8,14.168633) (12,14.661439) (16,9.317081) (24,8.704152)};\addlegendentry{40\%}
\end{axis}\end{tikzpicture}
\caption{移动目标的样本量与污染敏感性。线段仅连接离散实验格点，不表示格点间的理论插值。}
\label{fig:sensitivity-moving_outliers}
\end{figure}

\begin{table}[H]\centering\small
\caption{移动目标：模型选择及相对恒速基线的配对均值差。正差值表示所选方法误差较低。}

\begin{tabular}{lrrrrr}
\toprule
观测数 & 名义污染 & 实际污染 & 选择恒速 & CV可用 & 误差减少/m\\
\midrule
8 & 0\% & 0.0\% & 0.0\% & 0.0\% & -7.328\\
8 & 10\% & 12.5\% & 0.0\% & 0.0\% & 0.496\\
8 & 25\% & 25.0\% & 0.0\% & 0.0\% & 4.017\\
8 & 40\% & 37.5\% & 0.0\% & 0.0\% & 4.096\\
12 & 0\% & 0.0\% & 83.3\% & 100.0\% & -1.647\\
12 & 10\% & 8.3\% & 75.0\% & 100.0\% & -1.341\\
12 & 25\% & 25.0\% & 50.0\% & 100.0\% & -3.113\\
12 & 40\% & 41.7\% & 33.3\% & 100.0\% & 3.313\\
16 & 0\% & 0.0\% & 83.3\% & 100.0\% & -1.208\\
16 & 10\% & 12.5\% & 66.7\% & 100.0\% & -2.377\\
16 & 25\% & 25.0\% & 8.3\% & 100.0\% & -1.925\\
16 & 40\% & 37.5\% & 25.0\% & 100.0\% & -2.142\\
24 & 0\% & 0.0\% & 100.0\% & 100.0\% & -0.000\\
24 & 10\% & 8.3\% & 91.7\% & 100.0\% & -0.240\\
24 & 25\% & 25.0\% & 33.3\% & 100.0\% & -0.098\\
24 & 40\% & 41.7\% & 16.7\% & 100.0\% & -2.198\\
\bottomrule\end{tabular}\end{table}

离散观测数使实际污染比例与名义比例有所差异，表中分别列出。CV可用率为空表示该方法未输出此诊断；不能将缺失信息当作零。若冻结方法使用三折观测门控，则8条观测时部分训练折少于6条，按预定实现会禁止选择恒速；这项设计限制应与由污染或优化失败引起的选择错误区分。

\subsection{门控边界、污染交互与可追溯失败案例}
\label{sec:sensitivity-interpretation}

观测数量改变了信息量，也改变了模型可用性。三折时间交错划分要求每个训练折
至少有6份观测，因此计数条件为$n-\lceil n/3\rceil\geq6$。在$n=8$的96例中，
第一折训练集只有5份观测，程序直接停止CV，完成折数为0，运动可用标记全部为假，
最终全部输出静态模型。该运动无污染单元的均值误差为11.5264 m，恒速稳健
对照为4.1988 m。这里的静态选择来自实现门槛，不能解释成观测证据支持目标静止。

在$n\geq12$的288例中，三折均完成且运动模型全部可用。此时仍存在污染相关的
选择变化：真实运动案例在名义污染率$0,0.1,0.25,0.4$下，选择静态的次数依次为
$4/36,8/36,25/36,27/36$，对应平均位置误差为
3.6955、7.3217、8.8526、10.8942 m。CV比较历史观测上的标准化距离预测损失，
最终任务却评价参考时刻的位置；从距离误差到当前位置误差的放大受轨迹参数、
时间分布和几何条件影响，两类风险没有自动保持大小顺序的保证。

\begin{table}[H]
  \centering
  \small
  \caption{三个实际失败案例的配对位置误差。案例均保留于正式统计，单位为m。}
  \label{tab:sensitivity-failures}
  \begin{tabular}{llrrr}
    \toprule
    案例 & 最终选择 & 所选方法 & 恒速稳健 & 静态WLS\\
    \midrule
    A：运动，$n=16,\epsilon=0.1$ & 恒速 & 64.2834 & 64.2834 & 6.1442\\
    B：运动，$n=16,\epsilon=0.25$ & 静态 & 23.0179 & 3.1658 & 22.8234\\
    C：静态，$n=12,\epsilon=0.4$ & 恒速 & 41.5503 & 41.5503 & 9.3065\\
    \bottomrule
  \end{tabular}
\end{table}

案例A为本次最大误差，求解器仍报告成功、Jacobian秩为4、没有活动数值边界；
最小奇异值约0.16293，条件数约128.22。CV允许并选择恒速，但配对恒速基线也有
相同误差，说明失败不能只归因于模型选择。满秩与优化终止不能保证估计落在真解
附近，也不能代替理论局部误差界的邻域和噪声条件。案例B中运动模型可用，但CV
静态／运动得分为3.20774／3.51454，静态选择比配对恒速结果多出19.8521 m误差。
案例C则把零速度目标解释成运动，直接展示了静态非劣目标仍未得到保证。

为便于复核，表中A、B、C依次对应下列完整\texttt{case\_id}；在
\nolinkurl{results/sensitivity/predictions.json}中按该字段查找即可获得三种方法的
预测和CV诊断，在同目录\texttt{inputs.json}中可核对生成输入：
\begin{description}[leftmargin=1.6em,labelsep=0.4em,font=\normalfont\bfseries]
  \item[A] \nolinkurl{sensitivity-moving_outliers_n16_c1000-r05-921984553727470929}
  \item[B] \nolinkurl{sensitivity-moving_outliers_n16_c2500-r12-13744984418704934917}
  \item[C] \nolinkurl{sensitivity-static_outliers_n12_c4000-r04-14729045951357002381}
\end{description}

本分析为32格、每格12例的探索性描述。不同网格独立抽样，污染数量还经过
$\operatorname{round}(n\epsilon)$取整；观察到的交互和极端案例不能单独识别因果
机制，也不能精确定位污染崩溃阈值。上述解释结合代码规则与保存的诊断提出，
没有据此修改冻结方法、作逐格显著性宣称或建立新的覆盖保证。

\clearpage
\section{从已暴露失败到第二轮研究}
\label{sec:revision-design}

第一轮提供了一个真实但尚不充分的研究闭环。它发现的三个缺口直接决定第二轮问题：稀疏观测的三折训练集不足六条，导致运动模型被结构性排除；静态均值点估计的10\%容忍条件只在开发验证中满足，未在首次测试保持；总体边际区域掩盖运动异常子群的较低覆盖。第二轮因此研究折数、选择门槛及稳健组件的组合，并比较两种预先声明的区域校准方式。修订不以消除所有失败作为交付条件，最终结论服从新数据。

\subsection{旧结果是先验，新测试才是验证}
第二轮把第一轮结果作为已经暴露的先验，使用独立运行目录\path{runs/revision_reference}。发现集和验证集均包含四个常规场景及稀疏场景，种子分别为1646191与1657383，每场景分别20例与30例。第一轮数值源码、轨迹与冻结记录保持原字节；修订工具通过子类和独立入口扩展，实现与历史证据分开核验。

第二轮的全新主测试包含720例，压力测试400例，观测数边界测试420例，独立校准800例，覆盖检查1200例。机器可读方案\path{research/revision_plan.json}及评价器源码在新研究初始化时一并记录摘要。新案例只在方法冻结、源码和轨迹核验通过之后生成。因此，第二轮不能被描述成对第一轮研究路线的独立重复；它是\emph{依据已知失败作修订，再用新数据检验}的后续研究。

\subsection{更严格的工具约束}
第二轮显式开放原求解器已有的\texttt{cv\_folds}字段，让模型能够把折数写入可复核的方法定义。候选、单组件消融、诊断后修订、独立数学检查和验证证据分别留痕；冻结前要求候选与第一轮冻结方法及预设基线具有验证记录。工具也修订了第一轮消融说明继承父文本的问题：新登记同时保留父定义、实际改动和最终执行配置，避免把父算法的数学文字当成子算法的准确说明。

这些约束增加了过程可审查性，但并不证明规划器优于其他研究策略。两轮研究目标、开发数据和搜索空间不同；两次轨迹不能合并成同条件的两次随机重复，也不构成模型间的等预算对照。

\subsection{观测数与折数的结构关系}
若$K$个时间交错验证折的大小相差至多一条，最小训练折有
\[
 n_{\mathrm{train,min}}=n-\left\lceil\frac{n}{K}\right\rceil
\]
条观测。原实现要求运动训练折至少六条，因此必要计数条件为$n-\lceil n/K\rceil\ge6$。$n=8$时，$K=3$产生五条训练观测，而$K=4$产生六条。该变化只消除了计数障碍；时间跨度、几何秩、求解边界和噪声仍决定模型能否可靠拟合。

增加折数还会缩小每个验证折并增加拟合次数。实际评分汇总全部留出观测，故不能仅凭单折变小就断言总体评分更受异常影响；折数同时改变训练样本和预测误差的依赖关系，其净影响需要验证。更多拟合也提高计算成本。为避免仅展示有利的$n=8$，预登记边界测试完整覆盖$n=6,7,8,9,10,11,12$，在静态与运动无污染场景各生成30例，并保留所有不利结果。

\subsection{按可观测选模结果分组校准}
第一轮的子群欠覆盖提示误差分布异质性。本轮使用已有的条件共形框架\cite{vovk2012conditional}，不将分组校准称为原创方法。记完整观测为$Z$，冻结估计器为$\widehat x(Z)$，正尺度为$g(Z)$，并定义可观测分组
\[
 G(Z)\in\{\text{静态输出},\text{恒速输出}\}.
\]
它由冻结求解器的实际选择给出，不读取真实运动类别、污染标签或当前位置。对组$j$，取校准分数$S_i=\|X_i-\widehat x(Z_i)\|_2/g(Z_i)$中满足$G(Z_i)=j$的$m_j$项，并令
\[
 k_j=\left\lceil(m_j+1)(1-\alpha)\right\rceil,\qquad
 q_j=\begin{cases}S^{(j)}_{(k_j)},&k_j\le m_j,\\+\infty,&k_j>m_j.\end{cases}
\]
新区域为$C(Z)=\{x:\|x-\widehat x(Z)\|_2\le q_{G(Z)}g(Z)\}$。空组也按无界区域处理，不能用零半径或\texttt{null}表示已经获得有效的有限区域。

在估计器、尺度和分组函数预先冻结，校准与新案例独立同分布的条件下，给定新案例所属组及校准组大小，组内分数的交换性给出秩覆盖界。无并列时，新分数在$m_j+1$项中的秩均匀；有并列时采用不严格包含保持保守性。因此
\[
 \Pr\{X_{\rm new}\in C(Z_{\rm new})\mid G(Z_{\rm new})=j,\,M_j=m_j\}\ge1-\alpha.
\]
这个概率仍对校准分数取平均，不能解释为固定本次校准集后的保证。实际输出的静态组也不是“真实静止组”：误选模型会把不同物理场景混在一起。压力分布、边界分布、按半径拒绝的子集和任意精细真值场景均不自动获得这一条件覆盖。

本轮在相同800例校准与1200例覆盖数据上，同时报告全局与分组方法的覆盖、半径、接受率及真实场景内结果。分组样本减少可能增大分位数不确定性，改善覆盖也可能以更大半径为代价；选择最漂亮的单一指标不足以评价区域是否有用。

\section{第二轮真实智能体轨迹}

\label{sec:revision-trace}

本次第二轮保存了12个研究动作，发起13次模型请求，其中13次获得成功响应，其中审阅调用1次；实验与消融调用6次，登记4个新候选或消融版本，历时28.3分钟。正式状态为冻结。本记录由真实在线模型调用产生，未标为回放。网络等待包含在历时中，不能将其解释为定位求解耗时。各行决策来自模型返回，观测来自实际工具执行。

\begingroup\small\setlength{\tabcolsep}{4pt}

\begin{longtable}{@{}p{8mm}p{21mm}p{\dimexpr\linewidth-29mm-16pt\relax}@{}}

\caption{第二轮的完整动作顺序。方法以R编号简写，原始标识保留在方法表及日志中。}\\

\toprule 步 & 动作 & 公开决策依据与工具观测\\\midrule\endfirsthead

\toprule 步 & 动作 & 公开决策依据与工具观测\\\midrule\endhead

1 & 读取先验 & 尚无来源观察；先读取数学实现、第二轮预登记计划及已公开的v2结果，明确后续候选设计与验证约束。 读取原问题、实现数学、预登记计划及已暴露第一轮结果。\\[3pt]

2 & 登记候选 & 已读来源表明，三折在n=8时存在仅5个训练观测的折。先登记仅将v2折数改为4的候选，检验解除计数限制后的收益与代价；准确性改善尚无本轮证据。 登记R1；仅将v2的交错观测验证从三折改为四折：预计sparse恒速选择率提高且平均定位误差降低，static\_\allowbreak{}clean相对v2平均误差增幅不超过10\%，moving\_\allowbreak{}outliers平均误差不增加；预计平均耗时上升。四折只解除n=8的训练数量障碍，不保证几何可用性或\\[3pt]

3 & 执行实验 & 已读取来源并登记四折候选，但尚无本轮实验结果。先做配对发现集实验，以检验折数变化及静态、运动和成本取舍，并为后续诊断修订提供依据。 发现集五场景均值/m：R1 4.222；第一轮冻结 5.181；静态WLS 7.273；恒速稳健 4.644\\[3pt]

4 & 失败诊断 & 发现集上四折相对v2降低了稀疏和运动异常平均误差，但静态正常误差仍约为WLS的1.54倍，总体耗时增加约46\%。先诊断模型选择与失败案例，再据此修订；现有汇总尚不能确定退化机制。 诊断R1；最差返回案例误差24.221 m\\[3pt]

5 & 登记候选 & 诊断显示稀疏场景已选恒速14/20例，但最坏案例仍因验证得分选择静态；运动异常最坏例存在极大标准化残差。登记仅改变残差域的修订，检验污染评分与定位表现。静态最坏四例未返回，相关机制尚未确认。 登记R2；诊断驱动修订：四折已使sparse恒速选择达到14/20，但返回的三个稀疏失败例与一个运动异常失败例均因预测评分选择静态；后者含约20的标准化残差。仅将range改为log\_\allowbreak{}range，检验对数压缩与中心化是否改善评分及拟合。预测相对range\_\allowbreak{}cv4\_\allowbreak{}rep\\[3pt]

6 & 执行实验 & 对数四折修订已登记但尚无性能证据。下一步用相同发现案例与距离四折、v2及静态WLS配对比较，检验残差修改的效果及静态退化，保留充分预算完成消融、数学检查和验证。 发现集五场景均值/m：R2 5.036；R1 4.222；第一轮冻结 5.181；静态WLS 7.273\\[3pt]

7 & 单组件消融 & 对数修订仅改善静态正常误差，运动异常和稀疏改善预期被否定；同次实验耗时增幅约4.2\%，符合成本预期。以表现更好的距离四折为父方法，单独提高运动选择门槛，检验静态退化与运动收益的取舍，并完成必要消融。 发现集五场景均值/m：R1 4.222；R3 4.598。实际改动：\{"gate\_\allowbreak{}margin": 0.2\}\\[3pt]

8 & 数学核验 & 门槛消融降低静态正常误差，但运动异常误差增加约17.1\%，否定不超过10\%的预期。先完成独立数学检查，明确可证明结论的条件，再继续开发与验证。 核验：all。数值见证不替代数学证明。\\[3pt]

9 & 单组件消融 & 对数残差和提高运动门槛均出现误差退化，暂保留四折距离方法。下一步单独减少最终拟合起点，检验计算成本能否下降且不明显损害各目标场景精度；保留至少两次验证预算。 发现集五场景均值/m：R1 4.222；R4 4.222。实际改动：\{"starts": 1\}\\[3pt]

10 & 执行实验 & 单起点消融在发现集保持各场景误差近似不变，同次平均耗时降低约26.2\%，但静态正常退化仍未解决。现在验证单起点候选、父方法、v2及WLS，检验精度与成本取舍是否延续。 验证集五场景均值/m：R4 4.113；R1 4.064；第一轮冻结 5.221；静态WLS 7.414\\[3pt]

11 & 执行实验 & 已有来源、消融及数学检查证据，但三个参考方法缺少本轮验证，尚不能冻结。停止新增调参，用一次验证补齐对照并评估单起点候选的精度与耗时取舍。 验证集五场景均值/m：R4 4.113；静态稳健 7.033；恒速稳健 4.523；BIC门控 4.397\\[3pt]

12 & 冻结方法 & 全部冻结前提已满足，最新验证补齐五个参考。选择单起点候选，明确接受稀疏尾部代价与相对v2更高耗时；停止调参并冻结。 冻结R4；后续开发接口关闭。\\[3pt]

\bottomrule\end{longtable}\endgroup

\subsection{登记组合及冻结选择}

\begin{table}[H]\centering\small
\caption{新方法的主要组成。非CV模型不使用K和门槛；完整声明见原始注册表。}

\begin{tabularx}{\linewidth}{@{}p{8mm}p{36mm}>{\raggedright\arraybackslash}Xp{19mm}@{}}
\toprule
编号 & 原始标识 & 主要设置 & 验证均值/m\\
\midrule
R1 & range\_\allowbreak{}cv4\_\allowbreak{}reported & cv\_\allowbreak{}select; K=4; margin=0.05; soft\_\allowbreak{}l1; range; starts=3 & 4.064\\
R2 & log\_\allowbreak{}cv4\_\allowbreak{}reported & cv\_\allowbreak{}select; K=4; margin=0.05; soft\_\allowbreak{}l1; log\_\allowbreak{}range; starts=3 & --\\
R3 & range\_\allowbreak{}cv4\_\allowbreak{}margin20 & cv\_\allowbreak{}select; K=4; margin=0.2; soft\_\allowbreak{}l1; range; starts=3 & --\\
R4 & range\_\allowbreak{}cv4\_\allowbreak{}single\_\allowbreak{}start & cv\_\allowbreak{}select; K=4; margin=0.05; soft\_\allowbreak{}l1; range; starts=1 & 4.113\\
\bottomrule\end{tabularx}\end{table}

最终冻结方法为R4，其完整配置如下。

\begin{table}[H]\centering\small
\caption{第二轮冻结配置，测试后不再调整。}

\begin{tabularx}{\linewidth}{@{}>{\raggedright\arraybackslash}X>{\raggedright\arraybackslash}X@{}}
\toprule
字段 & 值\\
\midrule
consensus\_\allowbreak{}trials & 0\\
cv\_\allowbreak{}folds & 4\\
f\_\allowbreak{}scale & 1.5\\
gate\_\allowbreak{}margin & 0.05\\
log\_\allowbreak{}bias\_\allowbreak{}correction & True\\
log\_\allowbreak{}sd & 0.1\\
loss & soft\_\allowbreak{}l1\\
model & cv\_\allowbreak{}select\\
residual & range\\
ridge & 0.05\\
scale & reported\\
starts & 1\\
trim\_\allowbreak{}fraction & 0.0\\
\bottomrule\end{tabularx}\end{table}

\subsection{审阅反馈与研究结论原始记录}

\noindent\textbf{审阅1。} 在3次成功实验后记录。以下按真实审阅字段整理公开摘要，较长字段标明节选；这不是私有思维链。

\noindent\textbf{范围检查：} 通过。

\noindent\textbf{指出的问题：}\begin{enumerate}[leftmargin=2em,nosep]
\item 尚缺static\_\allowbreak{}robust、dynamic\_\allowbreak{}robust、legacy\_\allowbreak{}bic的validation证据，当前不能冻结；实验预算剩3次。
\item 静态正常目标仍失败：单起点验证均值1.248米，为WLS的约1.156倍，超过1.10界限。
\item 单起点满足预设均值容差，但稀疏均值较三起点增加约3.85\%，最大误差19.54米对16.10米；不能声称精度无损。
\item 单起点验证总体耗时较父方法下降约19.6\%，但较v2增加约22.8\%；未证明等计算预算优势，工作量口径也不可直接比较。
\item 验证集已暴露；后续依其选择候选属于模型选择，重复相同验证案例不是独立复现。未见禁用数据泄漏证据。
\item 对数修订及严格门槛的部分预测已失败，应保留负结果；新组合登记不证明学术新颖性。
\item 反射反例否定局部满秩蕴含全局唯一；标量偏差结论不推广至完整定位器，保形结论不提供当前实证覆盖保证。\end{enumerate}

\noindent\textbf{下一步建议：} 停止新增调参，下一次validation同时评估range\_\allowbreak{}cv4\_\allowbreak{}single\_\allowbreak{}start、static\_\allowbreak{}robust、dynamic\_\allowbreak{}robust、legacy\_\allowbreak{}bic，补齐参考对照。记录已有均值预测通过及静态非劣诊断失败、稀疏尾部代价，再按明确的精度与耗时取舍冻结新候选；不读取或生成holdout/test/calibration。

\noindent\textbf{证据边界：}\begin{itemize}[leftmargin=2em,nosep]
\item 仅审阅所供记录，未核验实现、哈希、样本独立性或计时环境。
\item 每场景验证仅30例，缺少配对差值不确定性及重复计时；均值阈值通过不等于统计非劣证明。
\item 四折解除计数障碍，同时改变评分与训练划分，误差改善不能完全归因于计数机制。
\item 本次同模型独立上下文审阅不是数学证明或独立经验复现。\end{itemize}

完整公开审阅字段见\nolinkurl{runs/revision_reference/state.json}中的critiques，及对应\nolinkurl{model_calls/}响应记录。独立提示不等于不同模型。

\noindent\textbf{冻结时的公开结论（保留原文）：}

依据本轮discovery与validation冻结现有完整配置，不再调参。已完成至少两个新组合、诊断后修订、成功单字段消融、数学检查及候选与全部五参考的验证覆盖。验证宏平均误差为4.113米，v2为5.221米、legacy\_\allowbreak{}bic为4.397米、dynamic\_\allowbreak{}robust为4.523米；候选运动异常与稀疏均值为7.313与6.608米，低于v2的8.599与10.687米，但高于dynamic\_\allowbreak{}robust的5.419与3.984米。静态正常1.248米，超过WLS的1.079米之1.10倍，预登记静态目标失败。

相对三起点父方法，同次验证总体耗时下降约19.6\%，宏均值由4.064升至4.113米，稀疏均值增加约3.85\%，最大误差由16.10升至19.54米；接受该精度与耗时取舍，不声称精度无损。相对v2同次耗时增加约22.8\%；最新同次候选耗时0.01547秒，static\_\allowbreak{}robust、dynamic\_\allowbreak{}robust、legacy\_\allowbreak{}bic分别为0.00327、0.00499、0.00779秒，无普遍速度优势。保留对数修订及严格门槛的失败预测。四折解除n=8训练数量障碍，但评分与训练划分同时改变，改善不能单独归因于计数机制。每场景仅30个验证案例，缺少配对不确定性和受控重复计时；验证参与模型选择，重复案例不构成独立复现。

反射恒等式否定局部满秩保证全局唯一；标量偏差公式不推广至完整定位器；保形结论依赖冻结与交换性等假设，尚无实证覆盖保证。不声称科学新颖性、统计非劣或工业问题已解决。未读取或生成本轮holdout、test或calibration；成功冻结立即结束开发。

完整提示、响应、失败和用量见\nolinkurl{runs/revision_reference/model_calls/}；完整动作、工具输入输出及证据链见同目录\texttt{events.jsonl}。

\section{第二轮冻结研究与全新独立评价}

\label{sec:revision-results}

第一轮失败结果是第二轮的已知先验；以下是新智能体冻结方法后，按预登记种子首次生成的独立数据。它们不构成对第一轮研究过程的盲重复，也未用于再次调整方法。

本轮冻结方法为\texttt{range\_\allowbreak{}cv4\_\allowbreak{}single\_\allowbreak{}start}。主测试、压力、观测数边界、校准和独立覆盖样本量依次为720、400、420、800、1200。

\begin{table}[H]\centering\small

\caption{第二轮所有预定比较的平均定位误差与主测试成本。误差单位为米。}

\begin{tabular}{@{}lrrrr@{}}\toprule

方法 & 主测试 & 压力 & 边界 & 主测试工作量\\\midrule

静态WLS & 6.643 & 12.772 & 6.177 & 19.5\\

静态稳健 & 6.171 & 12.849 & 6.229 & 21.4\\

恒速稳健 & 4.537 & 9.210 & 4.830 & 28.1\\

BIC门控 & 3.996 & 9.189 & 4.082 & 49.5\\

第一轮冻结 & 4.064 & 10.591 & 4.786 & 71.7\\

第二轮冻结 & 3.920 & 9.430 & 4.361 & 68.7\\

改为三折门控 & 4.076 & 10.587 & 4.769 & 55.1\\

改为四折门控$^{*}$ & 3.920 & 9.430 & 4.361 & 68.7\\

固定恒速 & 4.578 & 9.297 & 4.747 & 8.4\\

固定静态 & 6.180 & 12.855 & 6.220 & 5.9\\

\bottomrule\end{tabular}\end{table}

旧基线工作量为求解器nfev，新组合方法为含CV等步骤的残差调用数，口径并不完全相同；其比值不能直接解释为加速比，耗时只代表本次机器实测。

星号表示有效参数完全重合，复用同一预测而不视为独立方法：改为四折门控$=$第二轮冻结。

\begin{table}[H]\centering\small

\caption{主测试配对平均误差减少量及95\%区间；正值有利于第二轮冻结方法。}

\begin{tabular}{@{}lrrr@{}}\toprule

参照方法 & 减少量/m & 95\%区间/m & 相对减少/\%\\\midrule

静态WLS & 2.723 & [2.399, 3.036] & 40.99\\

静态稳健 & 2.252 & [1.944, 2.553] & 36.49\\

恒速稳健 & 0.617 & [0.330, 0.920] & 13.61\\

BIC门控 & 0.076 & [-0.158, 0.314] & 1.91\\

第一轮冻结 & 0.145 & [-0.043, 0.334] & 3.56\\

\bottomrule\end{tabular}\end{table}

区间在同一根种子与场景内重抽配对案例，并等权汇总固定单元；未校正多重比较，也不包含智能体搜索或方法选择的不确定性，因此属于探索性证据。

静止干净场景相对WLS的平均误差比为1.299，95\%区间为[1.165, 1.446]。预定容忍界为1.10；点估计未满足该界，配对差额的95\%区间上界未低于零。两者应分开陈述，不作确认性非劣结论。

\begin{table}[H]\centering\small

\caption{全局与观测选模分组的Mondrian区域：独立同分布覆盖检验。}

\begin{tabular}{@{}lrrrr@{}}\toprule

校准规则 & 名义/\% & 实测覆盖/\% & 平均半径/m & 无界数\\\midrule

global & 90 & 88.33 & 7.690 & 0\\

global & 95 & 94.58 & 11.779 & 0\\

mondrian & 90 & 88.17 & 7.247 & 0\\

mondrian & 95 & 95.00 & 10.802 & 0\\

\bottomrule\end{tabular}\end{table}

Mondrian分组仅使用冻结定位器根据观测选择的静态或恒速标签；真实场景仅用于事后评价。其保证依赖校准与测试可交换，并对随机校准评分取平均；不保证给定本次校准集、真实场景、按半径拒绝后的子集或分布偏移案例的覆盖。校准组为空或秩超过组样本数时区域无界，明确保留该代价。

全部场景、观测数边界、逐根种子结果、尾部与成本、压力与拒绝覆盖均保存在\nolinkurl{results/revision_benchmark/summary.json}及配套逐例文件。三个测试根种子代表数据批次，不能解释为三个独立智能体研究过程。

\subsection{全场景与跨种子的结果}

\begin{table}[H]\centering\small
\caption{主测试的逐场景均值，单位m。}

\begin{tabularx}{\linewidth}{@{}>{\raggedright\arraybackslash}X>{\raggedright\arraybackslash}X>{\raggedright\arraybackslash}X>{\raggedright\arraybackslash}X>{\raggedright\arraybackslash}X>{\raggedright\arraybackslash}X@{}}
\toprule
场景 & 静态WLS & 恒速稳健 & BIC门控 & 第一轮冻结 & 第二轮冻结\\
\midrule
移动正常 & 10.635 & 2.579 & 2.878 & 3.390 & 3.585\\
移动异常 & 11.168 & 7.245 & 9.805 & 9.186 & 8.562\\
静态正常 & 1.093 & 2.379 & 1.253 & 1.402 & 1.419\\
静态异常 & 3.675 & 5.944 & 2.048 & 2.280 & 2.112\\
\bottomrule\end{tabularx}\end{table}

\begin{table}[H]\centering\small
\caption{压力测试的逐场景均值，单位m。}

\begin{tabularx}{\linewidth}{@{}>{\raggedright\arraybackslash}X>{\raggedright\arraybackslash}X>{\raggedright\arraybackslash}X>{\raggedright\arraybackslash}X>{\raggedright\arraybackslash}X>{\raggedright\arraybackslash}X@{}}
\toprule
场景 & 静态WLS & 恒速稳健 & BIC门控 & 第一轮冻结 & 第二轮冻结\\
\midrule
加速运动 & 11.291 & 4.705 & 4.839 & 5.296 & 5.018\\
标定漂移 & 14.389 & 12.453 & 12.375 & 12.553 & 13.020\\
高污染 & 11.526 & 12.386 & 11.896 & 12.304 & 11.457\\
狭窄几何 & 15.540 & 10.464 & 10.835 & 11.380 & 10.660\\
稀疏观测 & 11.114 & 6.045 & 6.000 & 11.422 & 6.991\\
\bottomrule\end{tabularx}\end{table}

\begin{table}[H]\centering\small
\caption{三个新主测试数据批次的总体均值。它们不是三次智能体研究。}

\begin{tabularx}{\linewidth}{@{}>{\raggedright\arraybackslash}X>{\raggedright\arraybackslash}X>{\raggedright\arraybackslash}X>{\raggedright\arraybackslash}X@{}}
\toprule
根种子 & 第一轮冻结/m & 第二轮冻结/m & 恒速稳健/m\\
\midrule
1814711 & 4.228 & 4.047 & 4.812\\
1826013 & 3.839 & 3.724 & 4.013\\
1837901 & 4.126 & 3.988 & 4.786\\
\bottomrule\end{tabularx}\end{table}

\subsection{观测数边界与计算成本}

\begin{figure}[H]\centering
\begin{tikzpicture}\begin{axis}[width=.92\linewidth,height=6.6cm,xlabel={观测数量},ylabel={平均位置误差（m）},xtick={6,7,8,9,10,11,12},ymin=0,grid=major,legend style={at={(.5,1.04)},anchor=south,legend columns=3}]\addplot[thick,color=navy,mark=square*] coordinates {(6,1.672230) (7,1.793390) (8,1.887981) (9,1.955519) (10,1.757617) (11,1.510640) (12,1.689420)};\addlegendentry{第一轮冻结}
\addplot[thick,color=teal,mark=*] coordinates {(6,1.672405) (7,1.793398) (8,2.745938) (9,1.921730) (10,1.366116) (11,1.595903) (12,1.577056)};\addlegendentry{第二轮冻结}
\addplot[thick,color=orange,mark=triangle*] coordinates {(6,7.641730) (7,6.841625) (8,4.583020) (9,4.062102) (10,3.976314) (11,3.964654) (12,3.249252)};\addlegendentry{恒速稳健}\end{axis}\end{tikzpicture}
\caption{静态正常的完整观测数边界；每格30例，连线只表示离散格点。}
\label{fig:revision-edge-static_clean}
\end{figure}

\begin{figure}[H]\centering
\begin{tikzpicture}\begin{axis}[width=.92\linewidth,height=6.6cm,xlabel={观测数量},ylabel={平均位置误差（m）},xtick={6,7,8,9,10,11,12},ymin=0,grid=major,legend style={at={(.5,1.04)},anchor=south,legend columns=3}]\addplot[thick,color=navy,mark=square*] coordinates {(6,9.798988) (7,10.724031) (8,10.788175) (9,7.003231) (10,6.140565) (11,4.586809) (12,5.697341)};\addlegendentry{第一轮冻结}
\addplot[thick,color=teal,mark=*] coordinates {(6,9.565939) (7,10.723940) (8,6.216563) (9,6.609238) (10,5.850167) (11,5.202763) (12,4.218404)};\addlegendentry{第二轮冻结}
\addplot[thick,color=orange,mark=triangle*] coordinates {(6,7.708461) (7,5.612531) (8,5.106972) (9,4.495001) (10,3.214075) (11,3.738599) (12,3.421666)};\addlegendentry{恒速稳健}\end{axis}\end{tikzpicture}
\caption{移动正常的完整观测数边界；每格30例，连线只表示离散格点。}
\label{fig:revision-edge-moving_clean}
\end{figure}

\begin{table}[H]\centering\small
\caption{边界误差表（m），同时保留静态和运动表现。}

\begin{tabularx}{\linewidth}{@{}>{\raggedright\arraybackslash}X>{\raggedright\arraybackslash}X>{\raggedright\arraybackslash}X>{\raggedright\arraybackslash}X>{\raggedright\arraybackslash}X@{}}
\toprule
观测数 & 旧静态场景 & 新静态场景 & 旧运动场景 & 新运动场景\\
\midrule
6 & 1.672 & 1.672 & 9.799 & 9.566\\
7 & 1.793 & 1.793 & 10.724 & 10.724\\
8 & 1.888 & 2.746 & 10.788 & 6.217\\
9 & 1.956 & 1.922 & 7.003 & 6.609\\
10 & 1.758 & 1.366 & 6.141 & 5.850\\
11 & 1.511 & 1.596 & 4.587 & 5.203\\
12 & 1.689 & 1.577 & 5.697 & 4.218\\
\bottomrule\end{tabularx}\end{table}

\begin{table}[H]\centering\small
\caption{新主测试的成本与尾部。不同实现的工作量口径不同，耗时为本机单次测量。}

\begin{tabularx}{\linewidth}{@{}>{\raggedright\arraybackslash}X>{\raggedright\arraybackslash}X>{\raggedright\arraybackslash}X>{\raggedright\arraybackslash}X>{\raggedright\arraybackslash}X>{\raggedright\arraybackslash}X@{}}
\toprule
方法 & P90/m & P95/m & 最大/m & 工作量 & 耗时/ms\\
\midrule
静态WLS & 14.789 & 17.598 & 24.744 & 19.5 & 3.70\\
恒速稳健 & 9.302 & 13.443 & 55.311 & 28.1 & 6.76\\
第一轮冻结 & 10.008 & 14.171 & 55.311 & 71.7 & 20.14\\
第二轮冻结 & 9.973 & 13.970 & 55.311 & 68.7 & 16.92\\
\bottomrule\end{tabularx}\end{table}

\subsection{覆盖区间、分组差异和分布偏移}

\begin{table}[H]\centering\small
\caption{有限样本秩与实际校准阈值。空组或秩超过样本数时明确保留无界区域。}

\begin{tabularx}{\linewidth}{@{}>{\raggedright\arraybackslash}X>{\raggedright\arraybackslash}X>{\raggedright\arraybackslash}X>{\raggedright\arraybackslash}X>{\raggedright\arraybackslash}X>{\raggedright\arraybackslash}X>{\raggedright\arraybackslash}X@{}}
\toprule
规则 & 名义/\% & 校准组 & 样本 & 秩 & 阈值q & 无界\\
\midrule
全局 & 90 & 全部 & 800 & 721 & 3.244 & 否\\
全局 & 95 & 全部 & 800 & 761 & 4.968 & 否\\
分组 & 90 & 输出静态 & 480 & 433 & 4.200 & 否\\
分组 & 90 & 输出恒速 & 320 & 289 & 2.131 & 否\\
分组 & 95 & 输出静态 & 480 & 457 & 6.370 & 否\\
分组 & 95 & 输出恒速 & 320 & 305 & 3.087 & 否\\
\bottomrule\end{tabularx}\end{table}

\begin{table}[H]\centering\small
\caption{90\%名义水平的IID总体、真实场景与观测选模分组覆盖。区间为Wilson95\%，覆盖和区间单位为\%。}

\begin{tabularx}{\linewidth}{@{}>{\raggedright\arraybackslash}X>{\raggedright\arraybackslash}X>{\raggedright\arraybackslash}X>{\raggedright\arraybackslash}X>{\raggedright\arraybackslash}X>{\raggedright\arraybackslash}X@{}}
\toprule
规则 & 评价组 & 样本 & 覆盖 & 区间 & 半径/m\\
\midrule
全局 & 总体 & 1200 & 88.33 & [86.39, 90.03] & 7.690\\
全局 & 移动正常 & 270 & 90.74 & [86.69, 93.65] & 9.523\\
全局 & 移动异常 & 304 & 64.80 & [59.28, 69.96] & 10.706\\
全局 & 静态正常 & 329 & 100.00 & [98.85, 100.00] & 4.927\\
全局 & 静态异常 & 297 & 97.31 & [94.78, 98.63] & 5.997\\
全局 & 输出静态 & 744 & 84.27 & [81.48, 86.71] & 5.550\\
全局 & 输出恒速 & 456 & 94.96 & [92.55, 96.62] & 11.181\\
分组 & 总体 & 1200 & 88.17 & [86.22, 89.87] & 7.247\\
分组 & 移动正常 & 270 & 89.26 & [85.00, 92.42] & 6.603\\
分组 & 移动异常 & 304 & 66.78 & [61.30, 71.83] & 9.749\\
分组 & 静态正常 & 329 & 99.39 & [97.81, 99.83] & 5.647\\
分组 & 静态异常 & 297 & 96.63 & [93.91, 98.16] & 7.045\\
分组 & 输出静态 & 744 & 88.04 & [85.51, 90.18] & 7.186\\
分组 & 输出恒速 & 456 & 88.38 & [85.11, 91.00] & 7.347\\
\bottomrule\end{tabularx}\end{table}

\begin{table}[H]\centering\small
\caption{95\%名义水平的IID总体、真实场景与观测选模分组覆盖。区间为Wilson95\%，覆盖和区间单位为\%。}

\begin{tabularx}{\linewidth}{@{}>{\raggedright\arraybackslash}X>{\raggedright\arraybackslash}X>{\raggedright\arraybackslash}X>{\raggedright\arraybackslash}X>{\raggedright\arraybackslash}X>{\raggedright\arraybackslash}X@{}}
\toprule
规则 & 评价组 & 样本 & 覆盖 & 区间 & 半径/m\\
\midrule
全局 & 总体 & 1200 & 94.58 & [93.15, 95.73] & 11.779\\
全局 & 移动正常 & 270 & 93.33 & [89.71, 95.74] & 14.586\\
全局 & 移动异常 & 304 & 85.53 & [81.13, 89.04] & 16.398\\
全局 & 静态正常 & 329 & 100.00 & [98.85, 100.00] & 7.547\\
全局 & 静态异常 & 297 & 98.99 & [97.07, 99.66] & 9.186\\
全局 & 输出静态 & 744 & 92.07 & [89.91, 93.80] & 8.501\\
全局 & 输出恒速 & 456 & 98.68 & [97.16, 99.40] & 17.126\\
分组 & 总体 & 1200 & 95.00 & [93.62, 96.10] & 10.802\\
分组 & 移动正常 & 270 & 94.81 & [91.49, 96.89] & 9.612\\
分组 & 移动异常 & 304 & 87.17 & [82.94, 90.47] & 14.497\\
分组 & 静态正常 & 329 & 100.00 & [98.85, 100.00] & 8.514\\
分组 & 静态异常 & 297 & 97.64 & [95.22, 98.85] & 10.635\\
分组 & 输出静态 & 744 & 95.16 & [93.37, 96.48] & 10.900\\
分组 & 输出恒速 & 456 & 94.74 & [92.29, 96.44] & 10.642\\
\bottomrule\end{tabularx}\end{table}

\begin{table}[H]\centering\small
\caption{偏移分布上的描述性覆盖与半径，不赋予名义覆盖保证。}

\begin{tabularx}{\linewidth}{@{}>{\raggedright\arraybackslash}X>{\raggedright\arraybackslash}X>{\raggedright\arraybackslash}X>{\raggedright\arraybackslash}X>{\raggedright\arraybackslash}X>{\raggedright\arraybackslash}X@{}}
\toprule
规则 & 名义/\% & 总体 & 覆盖/\% & 半径/m & 无界数\\
\midrule
全局 & 90 & 压力 & 81.25 & 20.476 & 0\\
全局 & 90 & 边界 & 86.67 & 9.581 & 0\\
全局 & 95 & 压力 & 90.25 & 31.363 & 0\\
全局 & 95 & 边界 & 93.33 & 14.675 & 0\\
分组 & 90 & 压力 & 75.25 & 16.637 & 0\\
分组 & 90 & 边界 & 90.95 & 9.449 & 0\\
分组 & 95 & 压力 & 93.25 & 24.538 & 0\\
分组 & 95 & 边界 & 95.71 & 14.124 & 0\\
\bottomrule\end{tabularx}\end{table}

Wilson区间描述固定本次校准规则后，IID测试抽样对覆盖概率估计的不确定性；它与对校准评分取平均的共形保证不同，也不是第二次独立校准实验。逐真实场景和多个水平的区间未校正多重比较。

\subsection{第二轮修订回答了什么，仍未回答什么}

第二轮新主测试的第一轮冻结方法均值为4.064米，修订方法为3.920米，差值0.145米的配对95\%区间为$[-0.043,0.334]$。该区间跨零，不能据此宣布常规混合总体上有清晰的改进证据。修订方法相对BIC的差值区间同样跨零。与静态WLS的40.99\%均值改善回答的是另一个基线问题，不能替代对更强方法的比较。

稀疏与压力分布提供更具体的结果：在$n=8$的运动边界格点中，误差从10.788米降至6.217米，改善4.572米的配对区间为$[2.635,6.654]$；完整压力总体改善1.161米，区间$[0.661,1.666]$；完整边界总体改善0.425米，区间$[0.156,0.705]$。这些是预登记总体或格点上的探索性区间，没有校正全部比较的多重性。$n=8$上恒速稳健基线的误差仍更低，且该观测数下静态场景反而由1.888米退化至2.746米，说明解除计数排除没有解决稳健运动模型选择的全部困难。$n=6,7$也仍受训练数量限制。

静态目标没有被修订消除。在新主测试的180个静态正常案例中，修订误差1.419米，WLS为1.093米，均值比为1.299，95\%区间$[1.165,1.446]$。不仅点估计超过1.10容忍线，区间也全部高于该线。边界静态总体的比值为1.212，区间$[1.117,1.327]$。因此不能将总体平均收益解释为“标准静态场景不劣化”；该要求仍是明确失败项。

成本结果也有条件。正式新主测试中，修订方法的本机平均耗时约16.92毫秒，第一轮冻结方法约20.14毫秒，BIC约11.19毫秒。减少完整拟合起点并没有删除新增的交叉验证工作；不同模块的工作量和耗时不严格等价，而且这些是单轮、非随机化执行顺序的机器测量，不是稳定加速比或等预算算法优越性证明。

\subsection{分组校准改善的是哪一种平衡}
在相同的800例校准和1200例独立覆盖案例上，90\%全局区域实际覆盖88.33\%，平均半径7.690米；分组区域覆盖88.17\%，平均半径7.247米。分组没有提高这次总体覆盖。它将实际输出静态组与恒速组的覆盖由84.27\%和94.96\%，调整到88.04\%和88.38\%，体现的是两个可观测组之间的再平衡。

95\%名义水平下，全局覆盖94.58\%、平均半径11.779米，分组覆盖95.00\%、平均半径10.802米。这个单次结果支持继续研究分组规则的效率，但不能推广为所有水平都更好，也不应在看过结果后只报告95\%水平。两种水平与两种规则均预先声明并完整保留。

真实运动异常组在90\%分组区域下的覆盖仅66.78\%，全局规则为64.80\%。这是因为“输出模型类别”与“真实污染运动场景”不同，混合场景内的异质性没有被两个可观测组完全解释。拒绝后的子集与偏移分布结果也只能按各自指标描述，不能从总体共形定理自动获得条件保证。

本次90\%分组区域的覆盖率Wilson95\%区间为$[86.22\%,89.87\%]$，其上端低于90\%。这提示该固定校准规则的覆盖可能偏低，但不构成对共形秩定理的反例：定理的概率还对随机校准分数取平均，而本区间只描述固定校准后独立测试抽样的不确定性。报告同时保留了组样本量、秩、阈值与逐例分数，以便检查这种差异。

\subsection{过程与计算结果的独立核对}
两轮事件、方法注册、原始模型响应、工具证据和冻结源码均通过完整历史链审计，未报告不一致。第二轮独立统计脚本从保存的输入真值与输出位置重新计算距离、分场景统计、配对区间、有限样本校准秩、覆盖和拒绝结果，不调用定位器、不生成新数据，也不重新选择方法。123279项一致性检查通过，覆盖15860次实际预测、155组配对bootstrap比较及12个覆盖总体。原始核查输出保存在\path{verification/revision_results_check.json}；它检验记录与数值一致性，不能替代对数学假设或实验外推的判断。

第一轮冻结后计算为6192次，第二轮为15860次，合计22052次。第二轮声明的10个方法名称对应9个不同有效配置，四折消融是所选配置的准确别名，复用预测且明确标记。别名没有被计作额外独立计算；这也避免用相同预测的不同名称夸大实验工作量。另提供的新观测接口演示不计入以上正式评测次数。

\clearpage
\section{进一步的问题：允许拒绝是否足以消除定位歧义？}
\label{sec:rejection-question}

应用接口增加了按半径拒绝的功能，但“可以拒绝”本身不是可靠性的证明。本节在命题\ref{prop:mirror}的两世界构造上推导一个精确限制，再把后续研究问题写成可检验的约束。这是对既有不可区分性论证的直接延伸，不主张文献首次性；它属于报告中的数学分析，不计入智能体工具已经执行的数学检查次数。

\begin{proposition}[镜像歧义下的接受率与错误率限制]
\label{prop:reject-mirror}
设两个可能世界的当前位置为$x_0,x_1$，距离$D=\|x_0-x_1\|>0$，且完整观测$Z$的分布完全相同。两个世界使用同一可测规则，以$Z$及独立于$Z$和世界的同分布内部随机性$U$为输入，输出接受指标$A\in\{0,1\}$和位置$\widehat x$（拒绝时可任意定义该位置）。令共同接受概率为$p=\Pr(A=1)>0$。对任意$0\le r<D/2$，有
\begin{align}
 \max_{j\in\{0,1\}}\Pr_j\{\|\widehat x-x_j\|>r\mid A=1\}&\ge\frac12,\label{eq:reject-conditional}\\
 \max_{j\in\{0,1\}}\Pr_j\{A=1,\ \|\widehat x-x_j\|>r\}&\ge\frac p2.\label{eq:reject-joint}
\end{align}
因此，若要求每个世界的“接受且错误”概率均不超过$\alpha<1/2$，必须满足$p\le2\alpha$，即拒绝率至少为$1-2\alpha$。这里拒绝既不记为成功，也不记为一次已接受输出的错误；式\eqref{eq:reject-joint}采用的正是这个联合事件口径。
\end{proposition}

\begin{proof}
在两个世界使用相同分布的$(Z,U)$，其中$U$是估计器随机性。由于$2r<D$，两个闭球$B(x_0,r)$与$B(x_1,r)$互不相交。因此，对每次接受的相同输出，两个世界中至多有一个满足误差不超过$r$，即
\[
 \mathbf1_{\{A=1,\|\widehat x-x_0\|>r\}}+
 \mathbf1_{\{A=1,\|\widehat x-x_1\|>r\}}
 \ge \mathbf1_{\{A=1\}}.
\]
对共同分布积分，两个联合错误概率之和至少为$p$，其最大值至少为$p/2$。再除以两个世界共同的正接受概率$p$，得到条件结论。最后由联合错误均不超过$\alpha$推出$p\le2\alpha$。
\end{proof}

上述$1/2$界在允许规则知道两个候选点的两点模型中可以达到：以概率$p$接受，接受时等概率返回$x_0$或$x_1$。对任一世界，其条件错误率恰为$1/2$，联合错误率恰为$p/2$。这只是证明常数无法普遍收紧的理论规则，不是部署建议；它也不依赖声称某个局部最小值更可信。当$p=0$时联合错误为零，但接受后的条件错误率没有定义，不能将全拒绝解释为条件准确率100\%。

该结论说明三件事。第一，无信息的歧义不能仅通过缩小接受集合变成较高的最坏世界条件准确率；两个世界始终看到相同的接受事件。第二，允许大量拒绝可以降低联合错误，却会降低服务率。例如$\alpha=0.1$时，必要拒绝率至少为80\%。第三，这不否定前文的共形边际保证：保证所采用的分布、随机性、区域宽度及拒绝计分方式不同，不能互相替换。若输出区域允许同时覆盖两个镜像点，就必须承受相应的几何宽度。

\subsection{可供后续研究的受约束风险问题}
要提出有内容的改善目标，需要先排除上述不可区分性，或明确其出现概率。这里的$P$应定义为完整观测$Z$与真实当前位置、运动参数的\emph{联合分布}，而不是只给观测分布一个场景名称。固定有限观测时间窗，令$\mathcal P_0$中的轨迹以概率一为$x(t)=x$，$\mathcal P_1$中的轨迹以概率一为$x(t)=x+tv$且$v_{\min}\le\|v\|\le v_{\max}$，其中$0<v_{\min}\le v_{\max}<\infty$。位置参数域、污染机制、上报误差、观测数与时间设计必须在比较前固定；以下风险还要求相应条件二阶矩有限。

一种可进一步明确的\emph{全局}几何约束是：给定允许的理想锚点--时间设计$\mathcal D$，在规定参数域上，无噪声距离映射$F_{\mathcal D}(x,v)=(\|x+t_i v-a_i\|)_i$满足
\[
 \|F_{\mathcal D}(x,v)-F_{\mathcal D}(x',v')\|_2
 \ge c\|x-x'\|_2,\qquad c>0.
\]
它比较参数域内所有允许的参数对，因而排除了位置不同但理想距离完全相同的镜像；单独限制某拟合点Jacobian的最小奇异值不能替代它。哪些设计满足该条件、噪声与污染后能保留怎样的可辨识度，仍需证明；本报告不声称现有合成场景已经满足这一条件，也不把它本身当作有噪声风险界。

对规则$(\widehat x,A)$定义接受率及接受后的均方风险
\[
 a_P=\Pr_P(A=1),\qquad
 R_P=\mathbb E_P[\|\widehat x-x\|^2\mid A=1].
\]
接受指标仍只能依赖观测与预先指定的内部随机性。先固定基线估计器；其风险$R_P^{\rm WLS\mid A}$与$R_P^{\rm ref\mid A}$也必须在\emph{同一个由待评规则决定的接受事件}上计算；不能让各方法各自挑选不同的容易案例后直接比较。设$C$为统一口径下单次定位与接受判定的计算成本。一个明确的后续问题是：在
\[
 \inf_{P\in\mathcal P_0\cup\mathcal P_1}a_P\ge a_0>0,
 \qquad
 \sup_{P\in\mathcal P_0\cup\mathcal P_1}\mathbb E_P[C]\le B
\]
时，是否存在规则同时满足
\begin{align}
 \sup_{P\in\mathcal P_0}(R_P-R_P^{\rm WLS\mid A})&\le\delta_0,\\
 \sup_{P\in\mathcal P_1}(R_P-R_P^{\rm ref\mid A})&\le-\delta_1,
 \qquad \delta_1>0?
\end{align}
问题中的$(\delta_0,\delta_1,a_0,B)$给出静态容忍、移动收益、服务率和计算预算四项取舍。精度使用均方误差，因而不能把本报告的均值误差比较直接当作这些约束的证明。共同接受事件只保证比较同一批案例；若进一步声称等计算量优势，还须对基线施加同口径的预算。不同参数下可行集合可能为空：例如若移动分布类包含参照方法在接受事件上风险为零的分布，统一的正收益$\delta_1>0$立即不可能。命题\ref{prop:reject-mirror}提供另一类由不可区分性导致的障碍。

本次实验展示了促成该后续问题的具体失败机制，并实现了测量这些取舍的工具；它没有求出上述可行域，也没有给出统一的最优拒绝规则。下一步应先在明确的二维几何与污染分布类中求下界，再设计加入方向信息或轨迹先验的规则，最后进行预算匹配的数值检验。

\clearpage
\section{研究结论、局限与下一步问题}
\label{sec:discussion}

\subsection{可以确认的系统贡献}
本项目将大模型决策与数学计算分开实现：模型负责形成可检验的研究动作，求解器负责计算，独立评测负责检验。公开决策说明、实际模型原始输出、工具结果、阶段记忆和审阅意见可相互核对。这一结构使失败能够成为下一步输入，也使报告中的数值能够追溯到逐例计算。它体现了一个可运行科研智能体的开发与应用，而不是将预先写好的研究结论排成对话。

同时，系统不应被称为完全自主的科学家。问题选择、算子库、合成分布、工具接口、预算和统计协议由开发者设计；智能体在这个受限研究空间中组织假设、组合组件与修订方法。数学证明需要人工与符号层面的审查；独立critic调用仍来自同一配置模型，不能替代独立研究团队复现。两轮相继修订的真实轨迹也不足以证明这种规划策略优于随机、专家或其它模型；这需要额外的等预算、同问题、多次研究对照。

\subsection{数学与统计上的边界}
观测自加权偏差命题针对简化的标量对数正态模型。坐标误差、加性噪声、多锚点几何、截尾与局部优化同时出现时，其具体偏差并不由同一个标量公式决定。对数残差因此是一项受动机支持的候选，而非通用最优解。

镜像反例表明，数据本身的全局不可辨识性不能靠换损失函数消除。局部Jacobian满秩与求解器收敛，均不能排除另一处远离当前解的等价轨迹。局部误差界需要真值邻域、距离下界和曲率条件；它不能直接作为每个合成案例的保证书。

交叉验证衡量观测窗口内的预测适配，不保证选出使当前位置误差最小的运动类别。异常验证观测、非恒速运动和过少样本均可能破坏选择。共形区域的边际覆盖依赖完整案例之间的交换性，压力分布与拒绝后的条件子集超出了该保证。

\subsection{后续可检验的新问题}
本项目进一步提出一个适合继续研究的问题：在给定污染比例上界、最低几何信息量和有限求解预算下，能否构造允许拒绝的当前位置估计器，同时对静止情形给出非劣界、对移动情形给出误差改善，并控制拒绝率？这一问题必须明确分布类，否则镜像歧义会给出直接的不可能性障碍。

下一轮工作应首先处理可识别的部分：引入少量破除镜像的方向或地图信息，研究其最小数量；对非视距与设备标定漂移建立明确污染模型；用真正按设备或轨迹划分的现场数据检验区域覆盖；最后对不同智能体规划器开展多次独立重复与预算匹配的对照。它们是本项目留下的研究任务，而非本次已经完成的效果声明。

\clearpage
\bibliographystyle{unsrtnat}
\bibliography{../docs/references}
\clearpage
\appendix
\section{交付结构与复现实验}
\label{app:reproduction}

第一轮研究入口是 \path{run_agent.py}，独立评测入口为 \path{benchmark_v2_1.py}。第二轮对应入口为 \path{run_revision.py}与\path{benchmark_revision.py}，新版浏览器为\path{dashboard_revision.py}。数值求解器、基础研究工具与大模型运行时分别位于 \path{localization/}、\path{tools/} 和 \path{agent_runtime/}，第二轮扩展集中在\path{revision/}，从而可以分开检查数学算法与智能体决策。

\begin{table}[htbp]
\centering\small
\caption{项目文件与其用途。完整命令及版本条件见根目录README。}
\begin{tabularx}{\linewidth}{@{}p{49mm}Y@{}}
\toprule
路径 & 可核查内容\\
\midrule
\texttt{runs/live\_reference/} & 正式实际运行的状态、事件、模型调用、工具结果与冻结方法\\
\texttt{results/benchmark/} & 独立输入、逐例预测、统计摘要、消融、校准、协议与环境版本\\
\texttt{tests/} & 解析Jacobian、观测访问边界、交叉验证、控制器、证据与校准检验\\
\texttt{docs/} & 数学实现说明之外的研究背景、使用指南与文献数据库\\
\texttt{latex/} & 本报告的LaTeX主文件、章节与基于真实结果生成的表格和图形\\
\texttt{reporting/build\_report.py} & 从现有真实结果生成报告的数据部分，不重新选择定位方法\\
\bottomrule
\end{tabularx}
\end{table}

\subsection{三种不同含义的复现}
保存结果复查不产生新的模型请求。打开本地面板即可阅读方法提案、实验、失败和审阅意见，并核对逐例误差文件。确定性计算复现固定方法、输入和数值环境，再次运行求解与统计；壁钟时间和文件生成时刻可以变化，数值等价性应在指定容差内比较。

语言模型研究重跑具有不同含义。规划器受模型版本、服务状态与采样影响，新的调用可能作出不同决策；即使数据种子相同，也不能承诺复现同一条研究路径。应为新研究指定独立目录并保存新的模型原始输出。项目还提供显式的\texttt{replay}模式：它重新执行已保存的动作，测试控制器与工具的联接；它不会产生新的模型研究证据，面板和状态均将其标记为回放。

\subsection{基本命令}
在项目根目录创建Python环境并安装数值依赖后，可运行：
\begin{lstlisting}
python -m pip install -r requirements.txt
python -m unittest discover -s tests -v
python dashboard.py --run runs/live_reference --results results/benchmark
\end{lstlisting}
真实研究要求已安装Codex CLI并完成用户登录。登录信息由CLI自身管理，不写入项目。命令示例为：
\begin{lstlisting}
python run_agent.py --output runs/my_new_study
python benchmark_v2_1.py --run runs/my_new_study --output results/my_new_study
\end{lstlisting}
原\texttt{benchmark.py}保留为研究初始化版本及共享校验/统计库，不应直接作为正式评估入口。\texttt{research/evaluator\_amendment.md}记录冻结前、首次holdout前的入口修复与版本承诺。运行时不会在模型不可用时用规则伪造回应。调用失败会保留状态和错误证据；恢复必须保持研究目标、预算、种子与源代码一致。独立评测会验证冻结状态、方法注册表、验证集证据和源代码摘要，校验失败则停止。

报告使用XeLaTeX与BibTeX编译。Windows版采用系统宋体、黑体、Times New Roman与Consolas；新版在这些字体不可用时提供Fandol和TeX Gyre回退。所有科学图形使用PGFPlots或TikZ绘制，表格由保存的结果生成，无需图像生成模型。第二轮和完整PDF的命令见附录\ref{app:usability}以及根目录README。

\subsection{审计边界}
事件散列链、源代码摘要和不可覆盖的实验文件用于发现无意修改与版本不一致。它们不是受信第三方时间戳或密码学签名，不能抵御同时重写全部记录的攻击。模型服务未返回的信息不填入猜测值。公开发布前，提交者应补充自己的姓名、学号与课程信息，并检查模型调用记录中是否有不希望公开的研究文本；本交付不包含身份凭据或所提供课件原文。

\section{应用入口、交互核查与提交准备}
\label{app:usability}

\subsection{从保存研究到可使用的定位程序}
第二轮提供\path{locate_revision.py}。输入只含观测坐标\texttt{anchors}、相对时刻\texttt{times}、距离\texttt{ranges}与报告尺度\texttt{sigma}；位置单位为米，时刻单位为秒，参考时刻为0。输入不允许携带求解器可读取的目标真值或场景标签。文件\path{examples/observations.json}给出最小使用示例，其用途是接口演示，不属于工业现场验证。

\begin{lstlisting}
python locate_revision.py --input examples/observations.json \
  --run runs/revision_reference \
  --calibration results/revision_benchmark/conformal.json \
  --region mondrian --alpha 0.1 --max-radius 10 \
  --output results/my_location.json
\end{lstlisting}
上例用反斜杠示意分行；Windows PowerShell中请合并为一行运行，或依照README中的单行命令执行。程序验证冻结方法与校准记录的一致性，给出当前位置、速度、选中模型、几何诊断、名义区域及接受状态。选择\texttt{global}可使用全局校准，\texttt{mondrian}使用实际输出模型分组；两者均只能请求已保存的$\alpha$水平。

设置最大半径阈值后，半径超过阈值时状态为拒绝，点估计仍保留供诊断。未设阈值时状态为未应用拒绝规则，不能解释成通过了可靠性检查。对于空校准组或秩超出校准数量的情形，半径明确为无界；程序不会用零、缺失值或临时经验常数替代。现场新输入没有真值，所以输出也不计算无法观察到的实际误差。

\subsection{只读研究浏览器}
运行\path{dashboard_revision.py}后，用浏览器打开本机8766端口。页面可切换两轮研究，并提供五个视图：概览、智能体迭代、失败案例、校准与拒绝、观测数边界。它读取保存记录，不调用模型或定位器；使用者不需要模型账号即可检查成果。

迭代视图将公开模型决策摘要、假设、工具输入输出与审阅反馈连接起来。案例视图可以按方法误差排序，选择同一个\texttt{case\_id}比较多种方法。几何图使用相同米制比例显示锚点、历史测距圆、估计位置和五秒速度向量；评测真值明确标为离线专用，并可隐藏。因为圆属于不同历史时刻，不能把它们的静态交点当作正确的当前位置解释。

校准视图同时展示样本量、覆盖率、半径和拒绝曲线；边界图覆盖全部预登记观测数量。未开始、执行中、失败与完成分别显示，缺失证据不会自动变成零误差或虚构指标。浏览器用于核查研究结果，命令行应用用于计算一组新观测，两者承担不同的工作。

\subsection{环境自检、测试与报告重建}
\path{doctor.py}检查Python、数值依赖和可选的CLI版本。检查结果只表示本地依赖状况，不表示远端模型额度、登录或服务可用。已有结果浏览、数值重算、真实模型研究和LaTeX编译需要的环境不同，应按需要准备，详见README与分项指南。

\begin{lstlisting}
python doctor.py --json
python -m unittest discover -s tests -v
python reporting/compile_report.py
\end{lstlisting}
最后一个命令从保存证据更新实证章节，再执行XeLaTeX、BibTeX及必要的重复编译。未找到编译器、证据未完成、引用缺失、字体缺字或内容溢出时，脚本会报错，不把旧PDF标为新结果。\texttt{--xelatex}与\texttt{--bibtex}可指定编译器路径；\texttt{--skip-refresh}适合只修改文字和身份信息后的排版。新模板在Windows字体不可用时提供Fandol和TeX Gyre回退，实际交付版在Windows上编译并检查。

项目提供Windows与Linux、Python 3.11与3.12的离线CI配置。该配置文件随源码交付，不代表本次已经执行了远端四种组合；本机通过记录与平台范围以\path{verification/}为准。测试中的模拟模型用于验证预算、恢复、冻结、校准和输入边界，不能计作实际研究调用。

\subsection{建议的答辩展示路径}
先说明原工业问题与本项目合成子问题的区别，再从第一轮的静态退化或稀疏失败进入第二轮假设。接着展示一条真实“诊断--修订--消融--验证”的证据路径，最后用新的主测试、完整观测数边界与覆盖/半径取舍回答假设是否得到支持。对于没有改善的指标，应说明它约束了方法的适用范围。

完整代码、两轮原始轨迹、生成器输入、逐例结果、LaTeX和PDF一同构成提交项目。学生需补充姓名、学号及课程信息，并按课程要求说明开发辅助工具的使用。个人GitHub仓库由提交者创建并上传，本项目没有代为发布；课程原始课件、账号凭据及本地虚拟环境不属于交付内容。

\end{document}
